         \section{Products}
    \setcounter{figure}{1}
    \setcounter{subfigure}{1}
    \label{d4e6ddcad4e2d9e383c4732da6858c66}
         \subsection{ Introduction and recap}
    \nopagebreak
%%            \label{m39383} $ \hspace{-5pt}\begin{array}{cccccccccccc}   \includegraphics[width=0.75cm]{col11306.imgs/summary_fullmarks.png} &   \end{array} $ \hspace{2 pt}\raisebox{-5 pt}{} {(section shortcode: MG10060 )} \par 
    \label{m39383*cid2}
            \subsubsection{ Introduction}
            \nopagebreak
      \label{m39383*id266780}In this chapter you will learn how to work with algebraic expressions. You will recap some of the work on factorisation and multiplying out expressions that you learnt in earlier grades. This work will then be extended upon for Grade 10.\par 
    \label{m39383*cid3}
            \subsubsection{ Recap of Earlier Work}
            \nopagebreak
      \label{m39383*id266794}The following should be familiar. Examples are given as reminders.\par 
      \label{m39383*uid1}
            \subsubsubsection{ Parts of an Expression}
            \nopagebreak
        \label{m39383*id267144}Mathematical expressions are just like sentences and their parts have special names. You should be familiar with the following names used to describe the parts of a mathematical expression.\par 
        \label{m39383*uid2}\nopagebreak\noindent{}
          \settowidth{\mymathboxwidth}{\begin{equation}
    \begin{array}{cc}\hfill a\ensuremath{\cdot}{x}^{k}+b\ensuremath{\cdot}x+{c}^{m}=0\\ \hfill d\ensuremath{\cdot}{y}^{p}+e\ensuremath{\cdot}y+f\leq 0\end{array}\tag{8.1}
      \end{equation}
    }
    \typeout{Columnwidth = \the\columnwidth}\typeout{math as usual width = \the\mymathboxwidth}
    \ifthenelse{\lengthtest{\mymathboxwidth < \columnwidth}}{% if the math fits, do it again, for real
    \begin{equation}
    \begin{array}{cc}\hfill a\ensuremath{\cdot}{x}^{k}+b\ensuremath{\cdot}x+{c}^{m}=0\\ \hfill d\ensuremath{\cdot}{y}^{p}+e\ensuremath{\cdot}y+f\leq 0\end{array}\tag{8.1}
      \end{equation}
    }{% else, if it doesn't fit
    \setlength{\mymathboxwidth}{\columnwidth}
      \addtolength{\mymathboxwidth}{-48pt}
    \par\vspace{12pt}\noindent\begin{minipage}{\columnwidth}
    \parbox[t]{\mymathboxwidth}{\large$
    a\ensuremath{\cdot}{x}^{k}+b\ensuremath{\cdot}x+{c}^{m}=0d\ensuremath{\cdot}{y}^{p}+e\ensuremath{\cdot}y+f\leq 0$}\hfill
    \parbox[t]{48pt}{\raggedleft 
    (8.1)}
    \end{minipage}\vspace{12pt}\par
    }% end of conditional for this bit of math
    \typeout{math as usual width = \the\mymathboxwidth}
    % \textbf{m39383*uid3}\par
          \begin{table}[H]
    % \begin{table}[H]
    % \\ '' '0'
        \begin{center}
      \label{m39383*uid3}
    \noindent
    \tabletail{%
        \hline
        \multicolumn{2}{|p{\mytableboxwidth}|}{\raggedleft \small \sl continued on next page}\\
        \hline
      }
      \tablelasttail{}
      \begin{xtabular}[t]{|l|l|}\hline
        \textbf{Name} &
        \textbf{Examples (separated by commas)}% make-rowspan-placeholders
     \tabularnewline\cline{1-1}\cline{2-2}
      %--------------------------------------------------------------------
        term &
        $a\ensuremath{\cdot}{x}^{k}$ ,$b\ensuremath{\cdot}x$, ${c}^{m}$, $d\ensuremath{\cdot}{y}^{p}$, $e\ensuremath{\cdot}y$, $f$\hspace{20ex}% make-rowspan-placeholders
     \tabularnewline\cline{1-1}\cline{2-2}
      %--------------------------------------------------------------------
        expression &
        $a\ensuremath{\cdot}{x}^{k}+b\ensuremath{\cdot}x+{c}^{m}$, $d\ensuremath{\cdot}{y}^{p}+e\ensuremath{\cdot}y+f$% make-rowspan-placeholders
     \tabularnewline\cline{1-1}\cline{2-2}
      %--------------------------------------------------------------------
        coefficient &
        $a$, $b$, $d$, $e$% make-rowspan-placeholders
     \tabularnewline\cline{1-1}\cline{2-2}
      %--------------------------------------------------------------------
        exponent (or index) &
        $k$, $p$% make-rowspan-placeholders
     \tabularnewline\cline{1-1}\cline{2-2}
      %--------------------------------------------------------------------
        base &
        $x$, $y$, $c$% make-rowspan-placeholders
     \tabularnewline\cline{1-1}\cline{2-2}
      %--------------------------------------------------------------------
        constant &
        $a$, $b$, $c$, $d$, $e$, $f$% make-rowspan-placeholders
     \tabularnewline\cline{1-1}\cline{2-2}
      %--------------------------------------------------------------------
        variable &
        $x$, $y$% make-rowspan-placeholders
     \tabularnewline\cline{1-1}\cline{2-2}
      %--------------------------------------------------------------------
        equation &
                  $a\ensuremath{\cdot}{x}^{k}+b\ensuremath{\cdot}x+{c}^{m}=0$
                % make-rowspan-placeholders
     \tabularnewline\cline{1-1}\cline{2-2}
      %--------------------------------------------------------------------
        inequality &
                  $d\ensuremath{\cdot}{y}^{p}+e\ensuremath{\cdot}y+f\leq 0$
                % make-rowspan-placeholders
     \tabularnewline\cline{1-1}\cline{2-2}
      %--------------------------------------------------------------------
        binomial &
        expression with two terms% make-rowspan-placeholders
     \tabularnewline\cline{1-1}\cline{2-2}
      %--------------------------------------------------------------------
        trinomial &
        expression with three terms% make-rowspan-placeholders
     \tabularnewline\cline{1-1}\cline{2-2}
      %--------------------------------------------------------------------
    \end{xtabular}
      \end{center}
    \begin{center}{\small\bfseries Table 8.1}\end{center}
    \begin{caption}{\small\bfseries Table 8.1}\end{caption}
\end{table}
    \par
      \label{m39383*uid4}
            \subsubsubsection{ Product of Two Binomials}
            \nopagebreak
        \label{m39383*id268015}A \textsl{binomial} is a mathematical expression with two terms, e.g. $\left(ax+b\right)$ and $\left(cx+d\right)$. If these two binomials are multiplied, the following is the result:\par 
        \label{m39383*id268064}\nopagebreak\noindent{}\settowidth{\mymathboxwidth}{\begin{equation}
    \begin{array}{ccc}\hfill \left(a\ensuremath{\cdot}x+b\right)\left(c\ensuremath{\cdot}x+d\right)& =& \left(ax\right)\left(c\ensuremath{\cdot}x+d\right)+b\left(c\ensuremath{\cdot}x+d\right)\hfill \\ & =& \left(ax\right)\left(cx\right)+\left(ax\right)d+b\left(cx\right)+b\ensuremath{\cdot}d\hfill \\ & =& a{x}^{2}+x\left(ad+bc\right)+bd\hfill \end{array}\tag{8.2}
      \end{equation}
    }
    \typeout{Columnwidth = \the\columnwidth}\typeout{math as usual width = \the\mymathboxwidth}
    \ifthenelse{\lengthtest{\mymathboxwidth < \columnwidth}}{% if the math fits, do it again, for real
    \begin{equation}
    \begin{array}{ccc}\hfill \left(a\ensuremath{\cdot}x+b\right)\left(c\ensuremath{\cdot}x+d\right)& =& \left(ax\right)\left(c\ensuremath{\cdot}x+d\right)+b\left(c\ensuremath{\cdot}x+d\right)\hfill \\ & =& \left(ax\right)\left(cx\right)+\left(ax\right)d+b\left(cx\right)+b\ensuremath{\cdot}d\hfill \\ & =& a{x}^{2}+x\left(ad+bc\right)+bd\hfill \end{array}\tag{8.2}
      \end{equation}
    }{% else, if it doesn't fit
    \setlength{\mymathboxwidth}{\columnwidth}
      \addtolength{\mymathboxwidth}{-48pt}
    \par\vspace{12pt}\noindent\begin{minipage}{\columnwidth}
    \parbox[t]{\mymathboxwidth}{\large$
    \left(a\ensuremath{\cdot}x+b\right)\left(c\ensuremath{\cdot}x+d\right)=\left(ax\right)\left(c\ensuremath{\cdot}x+d\right)+b\left(c\ensuremath{\cdot}x+d\right)=\left(ax\right)\left(cx\right)+\left(ax\right)d+b\left(cx\right)+b\ensuremath{\cdot}d=a{x}^{2}+x\left(ad+bc\right)+bd$}\hfill
    \parbox[t]{48pt}{\raggedleft 
    (8.2)}
    \end{minipage}\vspace{12pt}\par
    }% end of conditional for this bit of math
    \typeout{math as usual width = \the\mymathboxwidth}
\par
            \label{m39383*secfhsst!!!underscore!!!id342}\vspace{.5cm} 
      \noindent
      \hspace*{-30pt}\includegraphics[width=0.5in]{col11306.imgs/pspencil2.png}   \raisebox{25mm}{   
      \begin{mdframed}[linewidth=4, leftmargin=40, rightmargin=40]  
      \begin{exercise}
    \noindent\textbf{Exercise 8.1: Product of two binomials }\label{m39383*probfhsst!!!underscore!!!id343}
        \label{m39383*id268279}Find the product of $\left(3x-2\right)\left(5x+8\right)$ \par 
        \vspace{5pt}
        \label{m39383*solfhsst!!!underscore!!!id346}\noindent\textbf{Solution to Exercise } \label{m39383*listfhsst!!!underscore!!!id346}\begin{enumerate}[noitemsep, label=\textbf{Step} \textbf{\arabic*}. ] 
            \leftskip=20pt\rightskip=\leftskip\item  
        \label{m39383*id268333}\nopagebreak\noindent{}
          \settowidth{\mymathboxwidth}{\begin{equation}
    \begin{array}{ccc}\hfill \left(3x-2\right)\left(5x+8\right)& =& \left(3x\right)\left(5x\right)+\left(3x\right)\left(8\right)+\left(-2\right)\left(5x\right)+\left(-2\right)\left(8\right)\hfill \\ & =& 15{x}^{2}+24x-10x-16\hfill \\ & =& 15{x}^{2}+14x-16\hfill \end{array}\tag{8.3}
      \end{equation}
    }
    \typeout{Columnwidth = \the\columnwidth}\typeout{math as usual width = \the\mymathboxwidth}
    \ifthenelse{\lengthtest{\mymathboxwidth < \columnwidth}}{% if the math fits, do it again, for real
    \begin{equation}
    \begin{array}{ccc}\hfill \left(3x-2\right)\left(5x+8\right)& =& \left(3x\right)\left(5x\right)+\left(3x\right)\left(8\right)+\left(-2\right)\left(5x\right)+\left(-2\right)\left(8\right)\hfill \\ & =& 15{x}^{2}+24x-10x-16\hfill \\ & =& 15{x}^{2}+14x-16\hfill \end{array}\tag{8.3}
      \end{equation}
    }{% else, if it doesn't fit
    \setlength{\mymathboxwidth}{\columnwidth}
      \addtolength{\mymathboxwidth}{-48pt}
    \par\vspace{12pt}\noindent\begin{minipage}{\columnwidth}
    \parbox[t]{\mymathboxwidth}{\large$
    \left(3x-2\right)\left(5x+8\right)=\left(3x\right)\left(5x\right)+\left(3x\right)\left(8\right)+\left(-2\right)\left(5x\right)+\left(-2\right)\left(8\right)=15{x}^{2}+24x-10x-16=15{x}^{2}+14x-16$}\hfill
    \parbox[t]{48pt}{\raggedleft 
    (8.3)}
    \end{minipage}\vspace{12pt}\par
    }% end of conditional for this bit of math
    \typeout{math as usual width = \the\mymathboxwidth}
        \end{enumerate}
    \end{exercise}
    \end{mdframed}
    }
    \noindent
        \label{m39383*id268534}The product of two identical binomials is known as the \textsl{square of the binomial} and is written as:\par 
        \label{m39383*id268543}\nopagebreak\noindent{}
          \settowidth{\mymathboxwidth}{\begin{equation}
    {\left(ax+b\right)}^{2}={a}^{2}{x}^{2}+2abx+{b}^{2}\tag{8.4}
      \end{equation}
    }
    \typeout{Columnwidth = \the\columnwidth}\typeout{math as usual width = \the\mymathboxwidth}
    \ifthenelse{\lengthtest{\mymathboxwidth < \columnwidth}}{% if the math fits, do it again, for real
    \begin{equation}
    {\left(ax+b\right)}^{2}={a}^{2}{x}^{2}+2abx+{b}^{2}\tag{8.4}
      \end{equation}
    }{% else, if it doesn't fit
    \setlength{\mymathboxwidth}{\columnwidth}
      \addtolength{\mymathboxwidth}{-48pt}
    \par\vspace{12pt}\noindent\begin{minipage}{\columnwidth}
    \parbox[t]{\mymathboxwidth}{\large$
    {\left(ax+b\right)}^{2}={a}^{2}{x}^{2}+2abx+{b}^{2}$}\hfill
    \parbox[t]{48pt}{\raggedleft 
    (8.4)}
    \end{minipage}\vspace{12pt}\par
    }% end of conditional for this bit of math
    \typeout{math as usual width = \the\mymathboxwidth}
        \label{m39383*id268608}If the two terms are $ax+b$\hspace{1ex} and $ax-b$\hspace{1ex} then their product is:\par 
        \label{m39383*id268642}\nopagebreak\noindent{}
          \settowidth{\mymathboxwidth}{\begin{equation}
    \left(ax+b\right)\left(ax-b\right)={a}^{2}{x}^{2}-{b}^{2}\tag{8.5}
      \end{equation}
    }
    \typeout{Columnwidth = \the\columnwidth}\typeout{math as usual width = \the\mymathboxwidth}
    \ifthenelse{\lengthtest{\mymathboxwidth < \columnwidth}}{% if the math fits, do it again, for real
    \begin{equation}
    \left(ax+b\right)\left(ax-b\right)={a}^{2}{x}^{2}-{b}^{2}\tag{8.5}
      \end{equation}
    }{% else, if it doesn't fit
    \setlength{\mymathboxwidth}{\columnwidth}
      \addtolength{\mymathboxwidth}{-48pt}
    \par\vspace{12pt}\noindent\begin{minipage}{\columnwidth}
    \parbox[t]{\mymathboxwidth}{\large$
    \left(ax+b\right)\left(ax-b\right)={a}^{2}{x}^{2}-{b}^{2}$}\hfill
    \parbox[t]{48pt}{\raggedleft 
    (8.5)}
    \end{minipage}\vspace{12pt}\par
    }% end of conditional for this bit of math
    \typeout{math as usual width = \the\mymathboxwidth}
        \label{m39383*id268705}This is known as the \textsl{difference of two squares}.\par 
      \label{m39383*uid5}
            \subsubsubsection{ Factorisation}
            \nopagebreak
        \label{m39383*id268725}Factorisation is the opposite of expanding brackets. For example expanding brackets would require $2\left(x+1\right)$ to be written as $2x+2$. Factorisation would be to start with $2x+2$\hspace{1ex} and to end up with $2\left(x+1\right)$. In previous grades, you factorised based on common factors and on difference of squares.\par 
        \label{m39383*uid6}
            \subsubsubsection{ Common Factors}
            \nopagebreak
          \label{m39383*id268808}Factorising based on common factors relies on there being common factors between your terms. For example, $2x-6{x}^{2}$\hspace{1ex}can be factorised as follows:\par 
          \label{m39383*id268835}\nopagebreak\noindent{}
            \settowidth{\mymathboxwidth}{\begin{equation}
    2x-6{x}^{2}=2x\left(1-3x\right)\tag{8.6}
      \end{equation}
    }
    \typeout{Columnwidth = \the\columnwidth}\typeout{math as usual width = \the\mymathboxwidth}
    \ifthenelse{\lengthtest{\mymathboxwidth < \columnwidth}}{% if the math fits, do it again, for real
    \begin{equation}
    2x-6{x}^{2}=2x\left(1-3x\right)\tag{8.6}
      \end{equation}
    }{% else, if it doesn't fit
    \setlength{\mymathboxwidth}{\columnwidth}
      \addtolength{\mymathboxwidth}{-48pt}
    \par\vspace{12pt}\noindent\begin{minipage}{\columnwidth}
    \parbox[t]{\mymathboxwidth}{\large$
    2x-6{x}^{2}=2x\left(1-3x\right)$}\hfill
    \parbox[t]{48pt}{\raggedleft 
    (8.6)}
    \end{minipage}\vspace{12pt}\par
    }% end of conditional for this bit of math
    \typeout{math as usual width = \the\mymathboxwidth}
\label{m39383*secfhsst!!!underscore!!!id565}
            \subsubsubsection{  Investigation : Common Factors }
            \nopagebreak
          \label{m39383*id268887}Find the highest common factors of the
following pairs of terms:\par 
    % \textbf{m39383*id268893}\par
          \begin{table}[H]
    % \begin{table}[H]
    % \\ 'id2889415' '1'
        \begin{center}
      \label{m39387*id273574}
    \noindent
    \tabletail{%
        \hline
        \multicolumn{2}{|p{\mytableboxwidth}|}{\raggedleft \small \sl continued on next page}\\
        \hline
      }
      \tablelasttail{}
      \begin{xtabular}[t]{|l|l|}\hline
        (a) $\left(-2{y}^{2}-4y+11\right)\left(5y-12\right)$ &
        (b) $\left(-11y+3\right)\left(-10{y}^{2}-7y-9\right)$% make-rowspan-placeholders
     \tabularnewline\cline{1-1}\cline{2-2}
      %--------------------------------------------------------------------
        (c) $\left(4{y}^{2}+12y+10\right)\left(-9{y}^{2}+8y+2\right)$ &
        (d) $\left(7{y}^{2}-6y-8\right)\left(-2y+2\right)$% make-rowspan-placeholders
     \tabularnewline\cline{1-1}\cline{2-2}
      %--------------------------------------------------------------------
        (e) $\left(10{y}^{5}+3\right)\left(-2{y}^{2}-11y+2\right)$ &
        (f) $\left(-12y-3\right)\left(12{y}^{2}-11y+3\right)$% make-rowspan-placeholders
     \tabularnewline\cline{1-1}\cline{2-2}
      %--------------------------------------------------------------------
        (g) $\left(-10\right)\left(2{y}^{2}+8y+3\right)$ &
        (h) $\left(2{y}^{6}+3{y}^{5}\right)\left(-5y-12\right)$% make-rowspan-placeholders
     \tabularnewline\cline{1-1}\cline{2-2}
      %--------------------------------------------------------------------
        (i) $\left(6{y}^{7}-8{y}^{2}+7\right)\left(-4y-3\right)\left(-6{y}^{2}-7y-11\right)$ &
        (j) $\left(-9{y}^{2}+11y+2\right)\left(8{y}^{2}+6y-7\right)$% make-rowspan-placeholders
     \tabularnewline\cline{1-1}\cline{2-2}
      %--------------------------------------------------------------------
        (k) $\left(8{y}^{5}+3{y}^{4}+2{y}^{3}\right)\left(5y+10\right)\left(12{y}^{2}+6y+6\right)$ &
        (l) $\left(-7y+11\right)\left(-12y+3\right)$% make-rowspan-placeholders
     \tabularnewline\cline{1-1}\cline{2-2}
      %--------------------------------------------------------------------
        (m) $\left(4{y}^{3}+5{y}^{2}-12y\right)\left(-12y-2\right)\left(7{y}^{2}-9y+12\right)$ &
        (n) $\left(7y+3\right)\left(7{y}^{2}+3y+10\right)$% make-rowspan-placeholders
     \tabularnewline\cline{1-1}\cline{2-2}
      %--------------------------------------------------------------------
        (o) $\left(9\right)\left(8{y}^{2}-2y+3\right)$ &
        (p) $\left(-12y+12\right)\left(4{y}^{2}-11y+11\right)$% make-rowspan-placeholders
     \tabularnewline\cline{1-1}\cline{2-2}
      %--------------------------------------------------------------------
        (q) $\left(-6{y}^{4}+11{y}^{2}+3y\right)\left(10y+4\right)\left(4y-4\right)$ &
        (r) $\left(-3{y}^{6}-6{y}^{3}\right)\left(11y-6\right)\left(10y-10\right)$% make-rowspan-placeholders
     \tabularnewline\cline{1-1}\cline{2-2}
      %--------------------------------------------------------------------
        (s) $\left(-11{y}^{5}+11{y}^{4}+11\right)\left(9{y}^{3}-7{y}^{2}-4y+6\right)$ &
        (t) $\left(-3y+8\right)\left(-4{y}^{3}+8{y}^{2}-2y+12\right)$% make-rowspan-placeholders
     \tabularnewline\cline{1-1}\cline{2-2}
      %--------------------------------------------------------------------
    \end{xtabular}
      \end{center}
    \begin{center}{\small\bfseries Table 8.5}\end{center}
    \begin{caption}{\small\bfseries Table 8.5}\end{caption}
\end{table}
    \par
  \newline
\newline
        \end{enumerate}
  \label{m39387**end}
\practiceinfo
\par 
 \par \begin{tabular}[h]{cccccc}
 (1.) 00iu& \end{tabular}
         \section{ Factorisation}
    \nopagebreak
%%            \label{m39394} $ \hspace{-5pt}\begin{array}{cccccccccccc}   \includegraphics[width=0.75cm]{col11306.imgs/summary_fullmarks.png} &   \includegraphics[width=0.75cm]{col11306.imgs/summary_video.png} &   \end{array} $ \hspace{2 pt}\raisebox{-5 pt}{} {(section shortcode: MG10063 )} \par 
           \subsubsubsection{ Factorisation}
            \nopagebreak
        \label{m39383*id268725}Factorisation is the opposite of expanding brackets. For example expanding brackets would require $2\left(x+1\right)$ to be written as $2x+2$. Factorisation would be to start with $2x+2$\hspace{1ex} and to end up with $2\left(x+1\right)$. In previous grades, you factorised based on common factors and on difference of squares.\par 
        \label{m39383*uid6}
            \subsubsubsection{ Common Factors}
            \nopagebreak
          \label{m39383*id268808}Factorising based on common factors relies on there being common factors between your terms. For example, $2x-6{x}^{2}$\hspace{1ex}can be factorised as follows:\par 
          \label{m39383*id268835}\nopagebreak\noindent{}
            \settowidth{\mymathboxwidth}{\begin{equation}
    2x-6{x}^{2}=2x\left(1-3x\right)\tag{8.6}
      \end{equation}
    }
    \typeout{Columnwidth = \the\columnwidth}\typeout{math as usual width = \the\mymathboxwidth}
    \ifthenelse{\lengthtest{\mymathboxwidth < \columnwidth}}{% if the math fits, do it again, for real
    \begin{equation}
    2x-6{x}^{2}=2x\left(1-3x\right)\tag{8.6}
      \end{equation}
    }{% else, if it doesn't fit
    \setlength{\mymathboxwidth}{\columnwidth}
      \addtolength{\mymathboxwidth}{-48pt}
    \par\vspace{12pt}\noindent\begin{minipage}{\columnwidth}
    \parbox[t]{\mymathboxwidth}{\large$
    2x-6{x}^{2}=2x\left(1-3x\right)$}\hfill
    \parbox[t]{48pt}{\raggedleft 
    (8.6)}
    \end{minipage}\vspace{12pt}\par
    }% end of conditional for this bit of math
    \typeout{math as usual width = \the\mymathboxwidth}
\label{m39383*secfhsst!!!underscore!!!id565}
            \subsubsubsection{  Investigation : Common Factors }
            \nopagebreak
          \label{m39383*id268887}Find the highest common factors of the
following pairs of terms:\par 
    % \textbf{m39383*id268893}\par
          \begin{table}[H]
    % \begin{table}[H]
    % \\ '' '0'
        \begin{center}
      \label{m39383*id268893}
    \noindent
    \tabletail{%
        \hline
        \multicolumn{5}{|p{\mytableboxwidth}|}{\raggedleft \small \sl continued on next page}\\
        \hline
      }
      \tablelasttail{}
      \begin{xtabular}[t]{|l|l|l|l|l|}\hline
        (a) $6y;18x$ &
        (b) $12mn;8n$ &
        (c) $3st;4su$ &
        (d) $18kl;9kp$ &
        (e) $abc;ac$% make-rowspan-placeholders
     \tabularnewline\cline{1-1}\cline{2-2}\cline{3-3}\cline{4-4}\cline{5-5}
      %--------------------------------------------------------------------
        (f) $2xy;4xyz$\hspace{5ex} &
        (g) $3uv;6u$ &
        (h) $9xy;15xz$\hspace{5ex} &
        (i) $24xyz;16yz$\hspace{5ex} &
        (j) $3m;45n$% make-rowspan-placeholders
     \tabularnewline\cline{1-1}\cline{2-2}\cline{3-3}\cline{4-4}\cline{5-5}
      %--------------------------------------------------------------------
    \end{xtabular}
      \end{center}
    \begin{center}{\small\bfseries Table 8.2}\end{center}
    \begin{caption}{\small\bfseries Table 8.2}\end{caption}
\end{table}
    \par
        \label{m39383*uid7}
            \subsubsubsection{ Difference of Two Squares}
            \nopagebreak
          \label{m39383*id269179}We have seen that:\par 
          \label{m39383*uid8}\nopagebreak\noindent{}
            \settowidth{\mymathboxwidth}{\begin{equation}
    \left(ax+b\right)\left(ax-b\right)={a}^{2}{x}^{2}-{b}^{2}\tag{8.7}
      \end{equation}
    }
    \typeout{Columnwidth = \the\columnwidth}\typeout{math as usual width = \the\mymathboxwidth}
    \ifthenelse{\lengthtest{\mymathboxwidth < \columnwidth}}{% if the math fits, do it again, for real
    \begin{equation}
    \left(ax+b\right)\left(ax-b\right)={a}^{2}{x}^{2}-{b}^{2}\tag{8.7}
      \end{equation}
    }{% else, if it doesn't fit
    \setlength{\mymathboxwidth}{\columnwidth}
      \addtolength{\mymathboxwidth}{-48pt}
    \par\vspace{12pt}\noindent\begin{minipage}{\columnwidth}
    \parbox[t]{\mymathboxwidth}{\large$
    \left(ax+b\right)\left(ax-b\right)={a}^{2}{x}^{2}-{b}^{2}$}\hfill
    \parbox[t]{48pt}{\raggedleft 
    (8.7)}
    \end{minipage}\vspace{12pt}\par
    }% end of conditional for this bit of math
    \typeout{math as usual width = \the\mymathboxwidth}
          \label{m39383*id269251}Since (8.7) is an equation, both sides are always equal. This means that an expression of the form:\par 
          \label{m39383*id269262}\nopagebreak\noindent{}
            \settowidth{\mymathboxwidth}{\begin{equation}
    {a}^{2}{x}^{2}-{b}^{2}\tag{8.8}
      \end{equation}
    }
    \typeout{Columnwidth = \the\columnwidth}\typeout{math as usual width = \the\mymathboxwidth}
    \ifthenelse{\lengthtest{\mymathboxwidth < \columnwidth}}{% if the math fits, do it again, for real
    \begin{equation}
    {a}^{2}{x}^{2}-{b}^{2}\tag{8.8}
      \end{equation}
    }{% else, if it doesn't fit
    \setlength{\mymathboxwidth}{\columnwidth}
      \addtolength{\mymathboxwidth}{-48pt}
    \par\vspace{12pt}\noindent\begin{minipage}{\columnwidth}
    \parbox[t]{\mymathboxwidth}{\large$
    {a}^{2}{x}^{2}-{b}^{2}$}\hfill
    \parbox[t]{48pt}{\raggedleft 
    (8.8)}
    \end{minipage}\vspace{12pt}\par
    }% end of conditional for this bit of math
    \typeout{math as usual width = \the\mymathboxwidth}
          \label{m39383*id269297}can be factorised to\par 
          \label{m39383*id269302}\nopagebreak\noindent{}
            \settowidth{\mymathboxwidth}{\begin{equation}
    \left(ax+b\right)\left(ax-b\right)\tag{8.9}
      \end{equation}
    }
    \typeout{Columnwidth = \the\columnwidth}\typeout{math as usual width = \the\mymathboxwidth}
    \ifthenelse{\lengthtest{\mymathboxwidth < \columnwidth}}{% if the math fits, do it again, for real
    \begin{equation}
    \left(ax+b\right)\left(ax-b\right)\tag{8.9}
      \end{equation}
    }{% else, if it doesn't fit
    \setlength{\mymathboxwidth}{\columnwidth}
      \addtolength{\mymathboxwidth}{-48pt}
    \par\vspace{12pt}\noindent\begin{minipage}{\columnwidth}
    \parbox[t]{\mymathboxwidth}{\large$
    \left(ax+b\right)\left(ax-b\right)$}\hfill
    \parbox[t]{48pt}{\raggedleft 
    (8.9)}
    \end{minipage}\vspace{12pt}\par
    }% end of conditional for this bit of math
    \typeout{math as usual width = \the\mymathboxwidth}
          \label{m39383*id269338}Therefore,\par 
          \label{m39383*id269343}\nopagebreak\noindent{}
            \settowidth{\mymathboxwidth}{\begin{equation}
    {a}^{2}{x}^{2}-{b}^{2}=\left(ax+b\right)\left(ax-b\right)\tag{8.10}
      \end{equation}
    }
    \typeout{Columnwidth = \the\columnwidth}\typeout{math as usual width = \the\mymathboxwidth}
    \ifthenelse{\lengthtest{\mymathboxwidth < \columnwidth}}{% if the math fits, do it again, for real
    \begin{equation}
    {a}^{2}{x}^{2}-{b}^{2}=\left(ax+b\right)\left(ax-b\right)\tag{8.10}
      \end{equation}
    }{% else, if it doesn't fit
    \setlength{\mymathboxwidth}{\columnwidth}
      \addtolength{\mymathboxwidth}{-48pt}
    \par\vspace{12pt}\noindent\begin{minipage}{\columnwidth}
    \parbox[t]{\mymathboxwidth}{\large$
    {a}^{2}{x}^{2}-{b}^{2}=\left(ax+b\right)\left(ax-b\right)$}\hfill
    \parbox[t]{48pt}{\raggedleft 
    (8.10)}
    \end{minipage}\vspace{12pt}\par
    }% end of conditional for this bit of math
    \typeout{math as usual width = \the\mymathboxwidth}
          \label{m39383*id269408}For example, ${x}^{2}-16$\hspace{1ex} can be written as $\left({x}^{2}-{4}^{2}\right)$ which is a difference of two squares. Therefore, the factors of ${x}^{2}-16$\hspace{1ex}are $\left(x-4\right)$ and $\left(x+4\right)$.\par 
\label{m39383*secfhsst!!!underscore!!!id708}\vspace{.5cm} 
      \noindent
      \hspace*{-30pt}\includegraphics[width=0.5in]{col11306.imgs/pspencil2.png}   \raisebox{25mm}{   
      \begin{mdframed}[linewidth=4, leftmargin=40, rightmargin=40]  
      \begin{exercise}
    \noindent\textbf{Exercise 8.2:  Factorisation }
          \label{m39383*probfhsst!!!underscore!!!id709}
          \label{m39383*id269529}Factorise completely: ${b}^{2}{y}^{5}-3ab{y}^{3}$ \par 
          \vspace{5pt}
          \label{m39383*solfhsst!!!underscore!!!id712}\noindent\textbf{Solution to Exercise } \label{m39383*listfhsst!!!underscore!!!id712}\begin{enumerate}[noitemsep, label=\textbf{Step} \textbf{\arabic*}. ] 
            \leftskip=20pt\rightskip=\leftskip\item  
          \label{m39383*id269591}\nopagebreak\noindent{}
            \settowidth{\mymathboxwidth}{\begin{equation}
    \begin{array}{ccc}\hfill {b}^{2}{y}^{5}-3ab{y}^{3}& =& b{y}^{3}\left(b{y}^{2}-3a\right)\hfill \end{array}\tag{8.11}
      \end{equation}
    }
    \typeout{Columnwidth = \the\columnwidth}\typeout{math as usual width = \the\mymathboxwidth}
    \ifthenelse{\lengthtest{\mymathboxwidth < \columnwidth}}{% if the math fits, do it again, for real
    \begin{equation}
    \begin{array}{ccc}\hfill {b}^{2}{y}^{5}-3ab{y}^{3}& =& b{y}^{3}\left(b{y}^{2}-3a\right)\hfill \end{array}\tag{8.11}
      \end{equation}
    }{% else, if it doesn't fit
    \setlength{\mymathboxwidth}{\columnwidth}
      \addtolength{\mymathboxwidth}{-48pt}
    \par\vspace{12pt}\noindent\begin{minipage}{\columnwidth}
    \parbox[t]{\mymathboxwidth}{\large$
    {b}^{2}{y}^{5}-3ab{y}^{3}=b{y}^{3}\left(b{y}^{2}-3a\right)$}\hfill
    \parbox[t]{48pt}{\raggedleft 
    (8.11)}
    \end{minipage}\vspace{12pt}\par
    }% end of conditional for this bit of math
    \typeout{math as usual width = \the\mymathboxwidth}
          \end{enumerate}
    \end{exercise}
    \end{mdframed}
    }
    \noindent
\par
            \label{m39383*secfhsst!!!underscore!!!id769}\vspace{.5cm} 
      \noindent
      \hspace*{-30pt}\includegraphics[width=0.5in]{col11306.imgs/pspencil2.png}   \raisebox{25mm}{   
      \begin{mdframed}[linewidth=4, leftmargin=40, rightmargin=40]  
      \begin{exercise}
    \noindent\textbf{Exercise 8.3:  Factorising binomials with a common bracket }
          \label{m39383*probfhsst!!!underscore!!!id770}
          \label{m39383*id269708}Factorise completely: $3a\left(a-4\right)-7\left(a-4\right)$ \par 
          \vspace{5pt}
          \label{m39383*solfhsst!!!underscore!!!id773}\noindent\textbf{Solution to Exercise } \label{m39383*listfhsst!!!underscore!!!id773}\begin{enumerate}[noitemsep, label=\textbf{Step} \textbf{\arabic*}. ] 
            \leftskip=20pt\rightskip=\leftskip\item  
          \newline
     $\left(a-4\right)$ is the common factor 
          \label{m39383*id269793}\nopagebreak\noindent{}
            \settowidth{\mymathboxwidth}{\begin{equation}
    \begin{array}{ccc}\hfill 3a\left(a-4\right)-7\left(a-4\right)& =& \left(a-4\right)\left(3a-7\right)\hfill \end{array}\tag{8.12}
      \end{equation}
    }
    \typeout{Columnwidth = \the\columnwidth}\typeout{math as usual width = \the\mymathboxwidth}
    \ifthenelse{\lengthtest{\mymathboxwidth < \columnwidth}}{% if the math fits, do it again, for real
    \begin{equation}
    \begin{array}{ccc}\hfill 3a\left(a-4\right)-7\left(a-4\right)& =& \left(a-4\right)\left(3a-7\right)\hfill \end{array}\tag{8.12}
      \end{equation}
    }{% else, if it doesn't fit
    \setlength{\mymathboxwidth}{\columnwidth}
      \addtolength{\mymathboxwidth}{-48pt}
    \par\vspace{12pt}\noindent\begin{minipage}{\columnwidth}
    \parbox[t]{\mymathboxwidth}{\large$
    3a\left(a-4\right)-7\left(a-4\right)=\left(a-4\right)\left(3a-7\right)$}\hfill
    \parbox[t]{48pt}{\raggedleft 
    (8.12)}
    \end{minipage}\vspace{12pt}\par
    }% end of conditional for this bit of math
    \typeout{math as usual width = \the\mymathboxwidth}
          \end{enumerate}
    \end{exercise}
    \end{mdframed}
    }
    \noindent
\par
            \label{m39383*secfhsst!!!underscore!!!id823}\vspace{.5cm} 
      \noindent
      \hspace*{-30pt}\includegraphics[width=0.5in]{col11306.imgs/pspencil2.png}   \raisebox{25mm}{   
      \begin{mdframed}[linewidth=4, leftmargin=40, rightmargin=40]  
      \begin{exercise}
    \noindent\textbf{Exercise 8.4:  Factorising using a switch around in brackets }
          \label{m39383*probfhsst!!!underscore!!!id824}
          \label{m39383*id269901}Factorise $5\left(a-2\right)-b\left(2-a\right)$ \par 
          \vspace{5pt}
          \label{m39383*solfhsst!!!underscore!!!id827}\noindent\textbf{Solution to Exercise } \label{m39383*listfhsst!!!underscore!!!id827}\begin{enumerate}[noitemsep, label=\textbf{Step} \textbf{\arabic*}. ] 
            \leftskip=20pt\rightskip=\leftskip\item  
          \label{m39383*id269994}\nopagebreak\noindent{}
            \settowidth{\mymathboxwidth}{\begin{equation}
    \begin{array}{ccc}\hfill 5\left(a-2\right)-b\left(2-a\right)& =& 5\left(a-2\right)-\left[-b\left(a-2\right)\right]\hfill \\ & =& 5\left(a-2\right)+b\left(a-2\right)\hfill \\ & =& \left(a-2\right)\left(5+b\right)\hfill \end{array}\tag{8.13}
      \end{equation}
    }
    \typeout{Columnwidth = \the\columnwidth}\typeout{math as usual width = \the\mymathboxwidth}
    \ifthenelse{\lengthtest{\mymathboxwidth < \columnwidth}}{% if the math fits, do it again, for real
    \begin{equation}
    \begin{array}{ccc}\hfill 5\left(a-2\right)-b\left(2-a\right)& =& 5\left(a-2\right)-\left[-b\left(a-2\right)\right]\hfill \\ & =& 5\left(a-2\right)+b\left(a-2\right)\hfill \\ & =& \left(a-2\right)\left(5+b\right)\hfill \end{array}\tag{8.13}
      \end{equation}
    }{% else, if it doesn't fit
    \setlength{\mymathboxwidth}{\columnwidth}
      \addtolength{\mymathboxwidth}{-48pt}
    \par\vspace{12pt}\noindent\begin{minipage}{\columnwidth}
    \parbox[t]{\mymathboxwidth}{\large$
    5\left(a-2\right)-b\left(2-a\right)=5\left(a-2\right)-\left[-b\left(a-2\right)\right]=5\left(a-2\right)+b\left(a-2\right)=\left(a-2\right)\left(5+b\right)$}\hfill
    \parbox[t]{48pt}{\raggedleft 
    (8.13)}
    \end{minipage}\vspace{12pt}\par
    }% end of conditional for this bit of math
    \typeout{math as usual width = \the\mymathboxwidth}
          \end{enumerate}
    \end{exercise}
    \end{mdframed}
    }
    \noindent
\label{m39383*secfhsst!!!underscore!!!id923}
            \subsubsubsection{  Recap }
            \nopagebreak
          \label{m39383*id270174}\begin{enumerate}[noitemsep, label=\textbf{\arabic*}. ] 
            \label{m39383*uid9}\item Find the products of:
    % \textbf{m39383*id270189}\par
          \begin{table}[H]
    % \begin{table}[H]
    % \\ 'id2885380' '1'
        \begin{center}
      \label{m39383*id270189}
    \noindent
    \tabletail{%
        \hline
        \multicolumn{3}{|p{\mytableboxwidth}|}{\raggedleft \small \sl continued on next page}\\
        \hline
      }
      \tablelasttail{}
      \begin{xtabular}[t]{|l|l|l|}\hline
        (a) $2y\left(y+4\right)$ &
        (b) $\left(y+5\right)\left(y+2\right)$ &
        (c) $\left(y+2\right)\left(2y+1\right)$% make-rowspan-placeholders
     \tabularnewline\cline{1-1}\cline{2-2}\cline{3-3}
      %--------------------------------------------------------------------
        (d) $\left(y+8\right)\left(y+4\right)$ &
        (e) $\left(2y+9\right)\left(3y+1\right)$ &
        (f) $\left(3y-2\right)\left(y+6\right)$% make-rowspan-placeholders
     \tabularnewline\cline{1-1}\cline{2-2}\cline{3-3}
      %--------------------------------------------------------------------
    \end{xtabular}
      \end{center}
    \begin{center}{\small\bfseries Table 8.3}\end{center}
    \begin{caption}{\small\bfseries Table 8.3}\end{caption}
\end{table}
    \par
  \newline
\newline
        \label{m39383*uid10}\item Factorise:
\label{m39383*id270414}\begin{enumerate}[noitemsep, label=\textbf{\alph*}. ] 
            \label{m39383*uid11}\item $2l+2w$\label{m39383*uid12}\item $12x+32y$\label{m39383*uid13}\item $6{x}^{2}+2x+10{x}^{3}$\label{m39383*uid14}\item $2x{y}^{2}+x{y}^{2}z+3xy$\label{m39383*uid15}\item $-2a{b}^{2}-4{a}^{2}b$\end{enumerate}
        \newline
\newline
        \label{m39383*uid16}\item Factorise completely:
    % \textbf{m39383*id270628}\par
          \begin{table}[H]
    % \begin{table}[H]
    % \\ 'id2885729' '1'
        \begin{center}
      \label{m39383*id270628}
    \noindent
    \tabletail{%
        \hline
        \multicolumn{3}{|p{\mytableboxwidth}|}{\raggedleft \small \sl continued on next page}\\
        \hline
      }
      \tablelasttail{}
      \begin{xtabular}[t]{|l|l|l|}\hline
        (a) $7a+4$ &
        (b) $20a-10$ &
        (c) $18ab-3bc$% make-rowspan-placeholders
     \tabularnewline\cline{1-1}\cline{2-2}\cline{3-3}
      %--------------------------------------------------------------------
        (d) $12kj+18kq$ &
        (e) $16{k}^{2}-4k$ &
        (f) $3{a}^{2}+6a-18$% make-rowspan-placeholders
     \tabularnewline\cline{1-1}\cline{2-2}\cline{3-3}
      %--------------------------------------------------------------------
        (g) $-6a-24$ &
        (h) $-2ab-8a$ &
        (i) $24kj-16{k}^{2}j$% make-rowspan-placeholders
     \tabularnewline\cline{1-1}\cline{2-2}\cline{3-3}
      %--------------------------------------------------------------------
        (j) $-{a}^{2}b-{b}^{2}a$ &
        (k) $12{k}^{2}j+24{k}^{2}{j}^{2}$ &
        (l) $72{b}^{2}q-18{b}^{3}{q}^{2}$% make-rowspan-placeholders
     \tabularnewline\cline{1-1}\cline{2-2}\cline{3-3}
      %--------------------------------------------------------------------
        (m) $4\left(y-3\right)+k\left(3-y\right)$ &
        (n) $a\left(a-1\right)-5\left(a-1\right)$ &
        (o) $bm\left(b+4\right)-6m\left(b+4\right)$% make-rowspan-placeholders
     \tabularnewline\cline{1-1}\cline{2-2}\cline{3-3}
      %--------------------------------------------------------------------
        (p) ${a}^{2}\left(a+7\right)+a\left(a+7\right)$ &
        (q) $3b\left(b-4\right)-7\left(4-b\right)$ &
        (r) ${a}^{2}{b}^{2}{c}^{2}-1$% make-rowspan-placeholders
     \tabularnewline\cline{1-1}\cline{2-2}\cline{3-3}
      %--------------------------------------------------------------------
    \end{xtabular}
      \end{center}
    \begin{center}{\small\bfseries Table 8.4}\end{center}
    \begin{caption}{\small\bfseries Table 8.4}\end{caption}
\end{table}
    \par
  \newline
\newline
        \end{enumerate}
  \label{m39383**end}
\practiceinfo
\par 
 \par \begin{tabular}[h]{cccccc}
 (1.) 00iv&  (2.) 00iw&  (3.) 00ix& \end{tabular}
    \label{m39394*cid5}
            \subsubsection{ Factorising a Quadratic}
            \nopagebreak
      \label{m39394*eip-218}
    \setcounter{subfigure}{0}
	\begin{figure}[H] % horizontal\label{m39394*factorisingquadratic}
    \textnormal{Khan Academy video on factorising a quadratic.}\vspace{.1in} \nopagebreak
  \label{m39394*yt-media2}\label{m39394*yt-video2}
            \raisebox{-5 pt}{ \includegraphics[width=0.5cm]{col11306.imgs/summary_www.png}} { (Video:  MG10064 )}
      \vspace{2pt}
    \vspace{.1in}
 \end{figure}       \par \label{m39394*eip-411}Factorisation can be seen as the reverse of calculating the product of factors. In order to factorise a quadratic, we need to find the factors which when multiplied together equal the original quadratic.\par 
      \label{m39394*id275057}Let us consider a quadratic that is of the form $a{x}^{2}+bx$\hspace{1ex}. We can see here that $x$ is a common factor of both terms. Therefore,\hspace{1ex}$a{x}^{2}+bx$\hspace{1ex}factorises to $x\left(ax+b\right)$. For example, $8{y}^{2}+4y$\hspace{1ex}factorises to\hspace{1ex}$4y\left(2y+1\right)$.\par 
      \label{m39394*id275188}Another type of quadratic is made up of the difference of squares. We know that:\par 
      \label{m39394*id275192}\nopagebreak\noindent{}
        \settowidth{\mymathboxwidth}{\begin{equation}
    \left(a+b\right)\left(a-b\right)={a}^{2}-{b}^{2}.\tag{8.19}
      \end{equation}
    }
    \typeout{Columnwidth = \the\columnwidth}\typeout{math as usual width = \the\mymathboxwidth}
    \ifthenelse{\lengthtest{\mymathboxwidth < \columnwidth}}{% if the math fits, do it again, for real
    \begin{equation}
    \left(a+b\right)\left(a-b\right)={a}^{2}-{b}^{2}.\tag{8.19}
      \end{equation}
    }{% else, if it doesn't fit
    \setlength{\mymathboxwidth}{\columnwidth}
      \addtolength{\mymathboxwidth}{-48pt}
    \par\vspace{12pt}\noindent\begin{minipage}{\columnwidth}
    \parbox[t]{\mymathboxwidth}{\large$
    \left(a+b\right)\left(a-b\right)={a}^{2}-{b}^{2}.$}\hfill
    \parbox[t]{48pt}{\raggedleft 
    (8.19)}
    \end{minipage}\vspace{12pt}\par
    }% end of conditional for this bit of math
    \typeout{math as usual width = \the\mymathboxwidth}
      \label{m39394*id275245}This is true for any values of $a$ and $b$, and more importantly since it is an equality, we can also write:\par 
      \label{m39394*id275270}\nopagebreak\noindent{}
        \settowidth{\mymathboxwidth}{\begin{equation}
    {a}^{2}-{b}^{2}=\left(a+b\right)\left(a-b\right).\tag{8.20}
      \end{equation}
    }
    \typeout{Columnwidth = \the\columnwidth}\typeout{math as usual width = \the\mymathboxwidth}
    \ifthenelse{\lengthtest{\mymathboxwidth < \columnwidth}}{% if the math fits, do it again, for real
    \begin{equation}
    {a}^{2}-{b}^{2}=\left(a+b\right)\left(a-b\right).\tag{8.20}
      \end{equation}
    }{% else, if it doesn't fit
    \setlength{\mymathboxwidth}{\columnwidth}
      \addtolength{\mymathboxwidth}{-48pt}
    \par\vspace{12pt}\noindent\begin{minipage}{\columnwidth}
    \parbox[t]{\mymathboxwidth}{\large$
    {a}^{2}-{b}^{2}=\left(a+b\right)\left(a-b\right).$}\hfill
    \parbox[t]{48pt}{\raggedleft 
    (8.20)}
    \end{minipage}\vspace{12pt}\par
    }% end of conditional for this bit of math
    \typeout{math as usual width = \the\mymathboxwidth}
      \label{m39394*id275324}This means that if we ever come across a quadratic that is made up of a difference of squares, we can immediately write down what the factors are.\par 
\label{m39394*secfhsst!!!underscore!!!id2008}\vspace{.5cm} 
      \noindent
      \hspace*{-30pt}\includegraphics[width=0.5in]{col11306.imgs/pspencil2.png}   \raisebox{25mm}{   
      \begin{mdframed}[linewidth=4, leftmargin=40, rightmargin=40]  
      \begin{exercise}
    \noindent\textbf{Exercise 8.7:  Difference of Squares }
      \label{m39394*probfhsst!!!underscore!!!id2009}
      \label{m39394*id275343}Find the factors of $9{x}^{2}-25$. \par 
      \vspace{5pt}
      \label{m39394*solfhsst!!!underscore!!!id2012}\noindent\textbf{Solution to Exercise } \label{m39394*listfhsst!!!underscore!!!id2012}\begin{enumerate}[noitemsep, label=\textbf{Step} \textbf{\arabic*}. ] 
            \leftskip=20pt\rightskip=\leftskip\item  
      \label{m39394*id275388}We see that the quadratic is a difference of squares because:\par 
      \label{m39394*id275392}\nopagebreak\noindent{}
        \settowidth{\mymathboxwidth}{\begin{equation}
    {\left(3x\right)}^{2}=9{x}^{2}\tag{8.21}
      \end{equation}
    }
    \typeout{Columnwidth = \the\columnwidth}\typeout{math as usual width = \the\mymathboxwidth}
    \ifthenelse{\lengthtest{\mymathboxwidth < \columnwidth}}{% if the math fits, do it again, for real
    \begin{equation}
    {\left(3x\right)}^{2}=9{x}^{2}\tag{8.21}
      \end{equation}
    }{% else, if it doesn't fit
    \setlength{\mymathboxwidth}{\columnwidth}
      \addtolength{\mymathboxwidth}{-48pt}
    \par\vspace{12pt}\noindent\begin{minipage}{\columnwidth}
    \parbox[t]{\mymathboxwidth}{\large$
    {\left(3x\right)}^{2}=9{x}^{2}$}\hfill
    \parbox[t]{48pt}{\raggedleft 
    (8.21)}
    \end{minipage}\vspace{12pt}\par
    }% end of conditional for this bit of math
    \typeout{math as usual width = \the\mymathboxwidth}
      \label{m39394*id275430}and\par 
      \label{m39394*id275435}\nopagebreak\noindent{}
        \settowidth{\mymathboxwidth}{\begin{equation}
    {5}^{2}=25.\tag{8.22}
      \end{equation}
    }
    \typeout{Columnwidth = \the\columnwidth}\typeout{math as usual width = \the\mymathboxwidth}
    \ifthenelse{\lengthtest{\mymathboxwidth < \columnwidth}}{% if the math fits, do it again, for real
    \begin{equation}
    {5}^{2}=25.\tag{8.22}
      \end{equation}
    }{% else, if it doesn't fit
    \setlength{\mymathboxwidth}{\columnwidth}
      \addtolength{\mymathboxwidth}{-48pt}
    \par\vspace{12pt}\noindent\begin{minipage}{\columnwidth}
    \parbox[t]{\mymathboxwidth}{\large$
    {5}^{2}=25.$}\hfill
    \parbox[t]{48pt}{\raggedleft 
    (8.22)}
    \end{minipage}\vspace{12pt}\par
    }% end of conditional for this bit of math
    \typeout{math as usual width = \the\mymathboxwidth}
      \item  
      \label{m39394*id275463}\nopagebreak\noindent{}
        \settowidth{\mymathboxwidth}{\begin{equation}
    9{x}^{2}-25={\left(3x\right)}^{2}-{5}^{2}\tag{8.23}
      \end{equation}
    }
    \typeout{Columnwidth = \the\columnwidth}\typeout{math as usual width = \the\mymathboxwidth}
    \ifthenelse{\lengthtest{\mymathboxwidth < \columnwidth}}{% if the math fits, do it again, for real
    \begin{equation}
    9{x}^{2}-25={\left(3x\right)}^{2}-{5}^{2}\tag{8.23}
      \end{equation}
    }{% else, if it doesn't fit
    \setlength{\mymathboxwidth}{\columnwidth}
      \addtolength{\mymathboxwidth}{-48pt}
    \par\vspace{12pt}\noindent\begin{minipage}{\columnwidth}
    \parbox[t]{\mymathboxwidth}{\large$
    9{x}^{2}-25={\left(3x\right)}^{2}-{5}^{2}$}\hfill
    \parbox[t]{48pt}{\raggedleft 
    (8.23)}
    \end{minipage}\vspace{12pt}\par
    }% end of conditional for this bit of math
    \typeout{math as usual width = \the\mymathboxwidth}
      \item  
      \label{m39394*id275517}\nopagebreak\noindent{}
        \settowidth{\mymathboxwidth}{\begin{equation}
    {\left(3x\right)}^{2}-{5}^{2}=\left(3x-5\right)\left(3x+5\right)\tag{8.24}
      \end{equation}
    }
    \typeout{Columnwidth = \the\columnwidth}\typeout{math as usual width = \the\mymathboxwidth}
    \ifthenelse{\lengthtest{\mymathboxwidth < \columnwidth}}{% if the math fits, do it again, for real
    \begin{equation}
    {\left(3x\right)}^{2}-{5}^{2}=\left(3x-5\right)\left(3x+5\right)\tag{8.24}
      \end{equation}
    }{% else, if it doesn't fit
    \setlength{\mymathboxwidth}{\columnwidth}
      \addtolength{\mymathboxwidth}{-48pt}
    \par\vspace{12pt}\noindent\begin{minipage}{\columnwidth}
    \parbox[t]{\mymathboxwidth}{\large$
    {\left(3x\right)}^{2}-{5}^{2}=\left(3x-5\right)\left(3x+5\right)$}\hfill
    \parbox[t]{48pt}{\raggedleft 
    (8.24)}
    \end{minipage}\vspace{12pt}\par
    }% end of conditional for this bit of math
    \typeout{math as usual width = \the\mymathboxwidth}
      \item  
      \label{m39394*id275585}The factors of $9{x}^{2}-25$\hspace{1ex} are $\left(3x-5\right)\left(3x+5\right)$. \par 
      \end{enumerate}
    \end{exercise}
    \end{mdframed}
    }
    \noindent
      \label{m39394*id275654}These types of quadratics are very simple to factorise. However, many quadratics do not fall into these categories and we need a more general method to factorise quadratics like ${x}^{2}-x-2$\hspace{1ex}?\par 
      \label{m39394*id275684}We can learn about how to factorise quadratics by looking at how two binomials are multiplied to get a quadratic. For example, $\left(x+2\right)\left(x+3\right)$ is multiplied out as:\par 
      \label{m39394*id275716}\nopagebreak\noindent{}
        \settowidth{\mymathboxwidth}{\begin{equation}
    \begin{array}{ccc}\hfill \left(x+2\right)\left(x+3\right)& =& x\left(x+3\right)+2\left(x+3\right)\hfill \\ & =& \left(x\right)\left(x\right)+3x+2x+\left(2\right)\left(3\right)\hfill \\ & =& {x}^{2}+5x+6.\hfill \end{array}\tag{8.25}
      \end{equation}
    }
    \typeout{Columnwidth = \the\columnwidth}\typeout{math as usual width = \the\mymathboxwidth}
    \ifthenelse{\lengthtest{\mymathboxwidth < \columnwidth}}{% if the math fits, do it again, for real
    \begin{equation}
    \begin{array}{ccc}\hfill \left(x+2\right)\left(x+3\right)& =& x\left(x+3\right)+2\left(x+3\right)\hfill \\ & =& \left(x\right)\left(x\right)+3x+2x+\left(2\right)\left(3\right)\hfill \\ & =& {x}^{2}+5x+6.\hfill \end{array}\tag{8.25}
      \end{equation}
    }{% else, if it doesn't fit
    \setlength{\mymathboxwidth}{\columnwidth}
      \addtolength{\mymathboxwidth}{-48pt}
    \par\vspace{12pt}\noindent\begin{minipage}{\columnwidth}
    \parbox[t]{\mymathboxwidth}{\large$
    \left(x+2\right)\left(x+3\right)=x\left(x+3\right)+2\left(x+3\right)=\left(x\right)\left(x\right)+3x+2x+\left(2\right)\left(3\right)={x}^{2}+5x+6.$}\hfill
    \parbox[t]{48pt}{\raggedleft 
    (8.25)}
    \end{minipage}\vspace{12pt}\par
    }% end of conditional for this bit of math
    \typeout{math as usual width = \the\mymathboxwidth}
      \label{m39394*id275871}We see that the ${x}^{2}$\hspace{1ex}term in the quadratic is the product of the $x$-terms in each bracket. Similarly, the 6 in the quadratic is the product of the 2 and 3 in the brackets. Finally, the middle term is the sum of two terms.\par 
      \label{m39394*id275901}So, how do we use this information to factorise the quadratic?\par 
      \label{m39394*id275905}Let us start with factorising ${x}^{2}+5x+6$ \hspace{1ex}and see if we can decide upon some general rules. Firstly, write down two brackets with an $x$ in each bracket and space for the remaining terms.\par 
      \label{m39394*id275944}\nopagebreak\noindent{}
        \settowidth{\mymathboxwidth}{\begin{equation}
    \left(x\phantom{\rule{2.em}{0ex}}\right)\left(x\phantom{\rule{2.em}{0ex}}\right)\tag{8.26}
      \end{equation}
    }
    \typeout{Columnwidth = \the\columnwidth}\typeout{math as usual width = \the\mymathboxwidth}
    \ifthenelse{\lengthtest{\mymathboxwidth < \columnwidth}}{% if the math fits, do it again, for real
    \begin{equation}
    \left(x\phantom{\rule{2.em}{0ex}}\right)\left(x\phantom{\rule{2.em}{0ex}}\right)\tag{8.26}
      \end{equation}
    }{% else, if it doesn't fit
    \setlength{\mymathboxwidth}{\columnwidth}
      \addtolength{\mymathboxwidth}{-48pt}
    \par\vspace{12pt}\noindent\begin{minipage}{\columnwidth}
    \parbox[t]{\mymathboxwidth}{\large$
    \left(x\phantom{\rule{2.em}{0ex}}\right)\left(x\phantom{\rule{2.em}{0ex}}\right)$}\hfill
    \parbox[t]{48pt}{\raggedleft 
    (8.26)}
    \end{minipage}\vspace{12pt}\par
    }% end of conditional for this bit of math
    \typeout{math as usual width = \the\mymathboxwidth}
      \label{m39394*id275980}Next, decide upon the factors of 6. Since the 6 is positive, these are:\par 
    % \textbf{m39394*id275986}\par
          \begin{table}[H]
    % \begin{table}[H]
    % \\ '' '0'
        \begin{center}
      \label{m39394*id275986}
    \noindent
    \tabletail{%
        \hline
        \multicolumn{2}{|p{\mytableboxwidth}|}{\raggedleft \small \sl continued on next page}\\
        \hline
      }
      \tablelasttail{}
      \begin{xtabular}[t]{|l|l|}\hline
    % My position: 0
    % my spanname: 
    % my ct of spanspec: 0
    % my column-count: 2
    \multicolumn{2}{|c|}{Factors of 6}
     \tabularnewline\cline{1-1}\cline{2-2}
      %--------------------------------------------------------------------
        1 &
        6% make-rowspan-placeholders
     \tabularnewline\cline{1-1}\cline{2-2}
      %--------------------------------------------------------------------
        2 &
        3% make-rowspan-placeholders
     \tabularnewline\cline{1-1}\cline{2-2}
      %--------------------------------------------------------------------
        -1 &
        -6% make-rowspan-placeholders
     \tabularnewline\cline{1-1}\cline{2-2}
      %--------------------------------------------------------------------
        -2 &
        -3% make-rowspan-placeholders
     \tabularnewline\cline{1-1}\cline{2-2}
      %--------------------------------------------------------------------
    \end{xtabular}
      \end{center}
    \begin{center}{\small\bfseries Table 8.6}\end{center}
    \begin{caption}{\small\bfseries Table 8.6}\end{caption}
\end{table}
    \par
      \label{m39394*id276096}Therefore, we have four possibilities:\par 
    % \textbf{m39394*id276099}\par
          \begin{table}[H]
    % \begin{table}[H]
    % \\ '' '0'
        \begin{center}
      \label{m39394*id276099}
    \noindent
    \tabletail{%
        \hline
        \multicolumn{4}{|p{\mytableboxwidth}|}{\raggedleft \small \sl continued on next page}\\
        \hline
      }
      \tablelasttail{}
      \begin{xtabular}[t]{|l|l|l|l|}\hline
        Option 1 &
        Option 2 &
        Option 3 &
        Option 4% make-rowspan-placeholders
     \tabularnewline\cline{1-1}\cline{2-2}\cline{3-3}\cline{4-4}
      %--------------------------------------------------------------------
                $\left(x+1\right)\left(x+6\right)$
               &
                $\left(x-1\right)\left(x-6\right)$
               &
                $\left(x+2\right)\left(x+3\right)$
               &
                $\left(x-2\right)\left(x-3\right)$
              % make-rowspan-placeholders
     \tabularnewline\cline{1-1}\cline{2-2}\cline{3-3}\cline{4-4}
      %--------------------------------------------------------------------
    \end{xtabular}
      \end{center}
    \begin{center}{\small\bfseries Table 8.7}\end{center}
    \begin{caption}{\small\bfseries Table 8.7}\end{caption}
\end{table}
    \par
      \label{m39394*id276261}Next, we expand each set of brackets to see which option gives us the correct middle term.\par 
    % \textbf{m39394*id276265}\par
          \begin{table}[H]
    % \begin{table}[H]
    % \\ '' '0'
        \begin{center}
      \label{m39394*id276265}
    \noindent
    \tabletail{%
        \hline
        \multicolumn{4}{|p{\mytableboxwidth}|}{\raggedleft \small \sl continued on next page}\\
        \hline
      }
      \tablelasttail{}
      \begin{xtabular}[t]{|l|l|l|l|}\hline
        Option 1 &
        Option 2 &
        Option 3 &
        Option 4% make-rowspan-placeholders
     \tabularnewline\cline{1-1}\cline{2-2}\cline{3-3}\cline{4-4}
      %--------------------------------------------------------------------
                $\left(x+1\right)\left(x+6\right)$
               &
                $\left(x-1\right)\left(x-6\right)$
               &
                $\left(x+2\right)\left(x+3\right)$
               &
                $\left(x-2\right)\left(x-3\right)$
              % make-rowspan-placeholders
     \tabularnewline\cline{1-1}\cline{2-2}\cline{3-3}\cline{4-4}
      %--------------------------------------------------------------------
                ${x}^{2}+7x+6$
               &
                ${x}^{2}-7x+6$
               &
                \uline{
                  ${x}^{2}+5x+6$
                }
               &
                ${x}^{2}-5x+6$
              % make-rowspan-placeholders
     \tabularnewline\cline{1-1}\cline{2-2}\cline{3-3}\cline{4-4}
      %--------------------------------------------------------------------
    \end{xtabular}
      \end{center}
    \begin{center}{\small\bfseries Table 8.8}\end{center}
    \begin{caption}{\small\bfseries Table 8.8}\end{caption}
\end{table}
    \par
      \label{m39394*id276547}We see that Option 3 (x+2)(x+3) is the correct solution. As you have seen that the process of factorising a quadratic is mostly trial and error, there is some information that can be used to simplify the process.\par 
      \label{m39394*uid20}
            \subsubsection{ Method: Factorising a Quadratic}
            \nopagebreak
        \label{m39394*id276561}\begin{enumerate}[noitemsep, label=\textbf{\arabic*}. ] 
            \label{m39394*uid21}\item First, divide the entire equation by any common factor of the coefficients so as to obtain an equation of the form $a{x}^{2}+bx+c=0$\hspace{1ex}where $a$, $b$ and $c$ have no common factors and $a$ is positive.
\label{m39394*uid22}\item Write down two brackets with an $x$ in each bracket and space for the remaining terms.
\label{m39394*uid23}\nopagebreak\noindent{}\settowidth{\mymathboxwidth}{\begin{equation}
    \left(x\phantom{\rule{2.em}{0ex}}\right)\left(x\phantom{\rule{2.em}{0ex}}\right)\tag{8.27}
      \end{equation}
    }
    \typeout{Columnwidth = \the\columnwidth}\typeout{math as usual width = \the\mymathboxwidth}
    \ifthenelse{\lengthtest{\mymathboxwidth < \columnwidth}}{% if the math fits, do it again, for real
    \begin{equation}
    \left(x\phantom{\rule{2.em}{0ex}}\right)\left(x\phantom{\rule{2.em}{0ex}}\right)\tag{8.27}
      \end{equation}
    }{% else, if it doesn't fit
    \setlength{\mymathboxwidth}{\columnwidth}
      \addtolength{\mymathboxwidth}{-48pt}
    \par\vspace{12pt}\noindent\begin{minipage}{\columnwidth}
    \parbox[t]{\mymathboxwidth}{\large$
    \left(x\phantom{\rule{2.em}{0ex}}\right)\left(x\phantom{\rule{2.em}{0ex}}\right)$}\hfill
    \parbox[t]{48pt}{\raggedleft 
    (8.27)}
    \end{minipage}\vspace{12pt}\par
    }% end of conditional for this bit of math
    \typeout{math as usual width = \the\mymathboxwidth}
    \label{m39394*uid24}\item Write down a set of factors for $a$ and $c$.
\label{m39394*uid25}\item Write down a set of options for the possible factors for the quadratic using the factors of $a$ and $c$.
\label{m39394*uid26}\item Expand all options to see which one gives you the correct answer.
\end{enumerate}
        \label{m39394*id276779}
          \textbf{There are some tips that you can keep in mind:}
        \par 
        \label{m39394*id276789}\begin{itemize}[noitemsep]
            \label{m39394*uid27}\item If $c$ is positive, then the factors of $c$ must be either both positive or both negative. The factors are both negative if $b$ is negative, and are both positive if $b$ is positive. If $c$ is negative, it means only one of the factors of $c$ is negative, the other one being positive.
\label{m39394*uid28}\item Once you get an answer, multiply out your brackets again just to make sure it really works.
\end{itemize}
\par
            \label{m39394*secfhsst!!!underscore!!!id2510}\vspace{.5cm} 
      \noindent
      \hspace*{-30pt}\includegraphics[width=0.5in]{col11306.imgs/pspencil2.png}   \raisebox{25mm}{   
      \begin{mdframed}[linewidth=4, leftmargin=40, rightmargin=40]  
      \begin{exercise}
    \noindent\textbf{Exercise 8.8:  Factorising a Quadratic }
        \label{m39394*probfhsst!!!underscore!!!id2511}
        \label{m39394*id276886}Find the factors of $3{x}^{2}+2x-1$. \par 
        \vspace{5pt}
        \label{m39394*solfhsst!!!underscore!!!id2514}\noindent\textbf{Solution to Exercise } \label{m39394*listfhsst!!!underscore!!!id2514}\begin{enumerate}[noitemsep, label=\textbf{Step} \textbf{\arabic*}. ] 
            \leftskip=20pt\rightskip=\leftskip\item  
        \label{m39394*id276979}The quadratic is in the required form.\par 
        \item  
        \label{m39394*uid29}\nopagebreak\noindent{}
          \settowidth{\mymathboxwidth}{\begin{equation}
    \left(x\phantom{\rule{2.em}{0ex}}\right)\left(x\phantom{\rule{2.em}{0ex}}\right)\tag{8.28}
      \end{equation}
    }
    \typeout{Columnwidth = \the\columnwidth}\typeout{math as usual width = \the\mymathboxwidth}
    \ifthenelse{\lengthtest{\mymathboxwidth < \columnwidth}}{% if the math fits, do it again, for real
    \begin{equation}
    \left(x\phantom{\rule{2.em}{0ex}}\right)\left(x\phantom{\rule{2.em}{0ex}}\right)\tag{8.28}
      \end{equation}
    }{% else, if it doesn't fit
    \setlength{\mymathboxwidth}{\columnwidth}
      \addtolength{\mymathboxwidth}{-48pt}
    \par\vspace{12pt}\noindent\begin{minipage}{\columnwidth}
    \parbox[t]{\mymathboxwidth}{\large$
    \left(x\phantom{\rule{2.em}{0ex}}\right)\left(x\phantom{\rule{2.em}{0ex}}\right)$}\hfill
    \parbox[t]{48pt}{\raggedleft 
    (8.28)}
    \end{minipage}\vspace{12pt}\par
    }% end of conditional for this bit of math
    \typeout{math as usual width = \the\mymathboxwidth}
        \label{m39394*id277035}Write down a set of factors for $a$ and $c$.
The possible factors for $a$ are: (1,3).
The possible factors for $c$ are: (-1,1) or (1,-1).\par 
        \label{m39394*id277075}Write down a set of options for the possible factors of the quadratic using the factors of $a$ and $c$.
Therefore, there are two possible options.\par 
    % \textbf{m39394*id277097}\par
          \begin{table}[H]
    % \begin{table}[H]
    % \\ 'id2892888' '1'
        \begin{center}
      \label{m39394*id277097}
    \noindent
    \tabletail{%
        \hline
        \multicolumn{2}{|p{\mytableboxwidth}|}{\raggedleft \small \sl continued on next page}\\
        \hline
      }
      \tablelasttail{}
      \begin{xtabular}[t]{|l|l|}\hline
        Option 1 &
        Option 2% make-rowspan-placeholders
     \tabularnewline\cline{1-1}\cline{2-2}
      %--------------------------------------------------------------------
                  $\left(x-1\right)\left(3x+1\right)$
                 &
                  $\left(x+1\right)\left(3x-1\right)$
                % make-rowspan-placeholders
     \tabularnewline\cline{1-1}\cline{2-2}
      %--------------------------------------------------------------------
                  $3{x}^{2}-2x-1$
                 &
                  \uline{
                    $3{x}^{2}+2x-1$
                  }
                % make-rowspan-placeholders
     \tabularnewline\cline{1-1}\cline{2-2}
      %--------------------------------------------------------------------
    \end{xtabular}
      \end{center}
    \begin{center}{\small\bfseries Table 8.9}\end{center}
    \begin{caption}{\small\bfseries Table 8.9}\end{caption}
\end{table}
    \par
        \item  
        \label{m39394*id277257}\nopagebreak\noindent{}
          \settowidth{\mymathboxwidth}{\begin{equation}
    \begin{array}{ccc}\hfill \left(x+1\right)\left(3x-1\right)& =& x\left(3x-1\right)+1\left(3x-1\right)\hfill \\ & =& \left(x\right)\left(3x\right)+\left(x\right)\left(-1\right)+\left(1\right)\left(3x\right)+\left(1\right)\left(-1\right)\hfill \\ & =& 3{x}^{2}-x+3x-1\hfill \\ & =& {x}^{2}+2x-1.\hfill \end{array}\tag{8.29}
      \end{equation}
    }
    \typeout{Columnwidth = \the\columnwidth}\typeout{math as usual width = \the\mymathboxwidth}
    \ifthenelse{\lengthtest{\mymathboxwidth < \columnwidth}}{% if the math fits, do it again, for real
    \begin{equation}
    \begin{array}{ccc}\hfill \left(x+1\right)\left(3x-1\right)& =& x\left(3x-1\right)+1\left(3x-1\right)\hfill \\ & =& \left(x\right)\left(3x\right)+\left(x\right)\left(-1\right)+\left(1\right)\left(3x\right)+\left(1\right)\left(-1\right)\hfill \\ & =& 3{x}^{2}-x+3x-1\hfill \\ & =& {x}^{2}+2x-1.\hfill \end{array}\tag{8.29}
      \end{equation}
    }{% else, if it doesn't fit
    \setlength{\mymathboxwidth}{\columnwidth}
      \addtolength{\mymathboxwidth}{-48pt}
    \par\vspace{12pt}\noindent\begin{minipage}{\columnwidth}
    \parbox[t]{\mymathboxwidth}{\large$
    \left(x+1\right)\left(3x-1\right)=x\left(3x-1\right)+1\left(3x-1\right)=\left(x\right)\left(3x\right)+\left(x\right)\left(-1\right)+\left(1\right)\left(3x\right)+\left(1\right)\left(-1\right)=3{x}^{2}-x+3x-1={x}^{2}+2x-1.$}\hfill
    \parbox[t]{48pt}{\raggedleft 
    (8.29)}
    \end{minipage}\vspace{12pt}\par
    }% end of conditional for this bit of math
    \typeout{math as usual width = \the\mymathboxwidth}
        \item  
        \label{m39394*id277481}The factors of $3{x}^{2}+2x-1$\hspace{1ex}are $\left(x+1\right)$ and $\left(3x-1\right)$.
 \par 
        \end{enumerate}
    \end{exercise}
    \end{mdframed}
    }
    \noindent
\label{m39394*secfhsst!!!underscore!!!id2756}
            \subsubsection{  Factorising a Trinomial }
            \nopagebreak
        \label{m39394*id277570}\begin{enumerate}[noitemsep, label=\textbf{\arabic*}. ] 
            \label{m39394*uid30}\item Factorise the following:
    % \textbf{m39394*id277586}\par
          \begin{table}[H]
    % \begin{table}[H]
    % \\ 'id2893661' '1'
        \begin{center}
      \label{m39394*id277586}
    \noindent
    \tabletail{%
        \hline
        \multicolumn{3}{|p{\mytableboxwidth}|}{\raggedleft \small \sl continued on next page}\\
        \hline
      }
      \tablelasttail{}
      \begin{xtabular}[t]{|l|l|l|}\hline
        (a) ${x}^{2}+8x+15$\hspace{1ex} &
        (b) ${x}^{2}+10x+24$\hspace{1ex} &
        (c) ${x}^{2}+9x+8$\hspace{1ex}% make-rowspan-placeholders
     \tabularnewline\cline{1-1}\cline{2-2}\cline{3-3}
      %--------------------------------------------------------------------
        (d) ${x}^{2}+9x+14$\hspace{1ex} &
        (e) ${x}^{2}+15x+36$\hspace{1ex} &
        (f) ${x}^{2}+12x+36$\hspace{1ex}% make-rowspan-placeholders
     \tabularnewline\cline{1-1}\cline{2-2}\cline{3-3}
      %--------------------------------------------------------------------
    \end{xtabular}
      \end{center}
    \begin{center}{\small\bfseries Table 8.10}\end{center}
    \begin{caption}{\small\bfseries Table 8.10}\end{caption}
\end{table}
    \par
          \label{m39394*uid31}\item Factorise the following:
\label{m39394*id277790}\begin{enumerate}[noitemsep, label=\textbf{\alph*}. ] 
            \label{m39394*uid32}\item ${x}^{2}-2x-15$\label{m39394*uid33}\item ${x}^{2}+2x-3$\label{m39394*uid34}\item ${x}^{2}+2x-8$\label{m39394*uid35}\item ${x}^{2}+x-20$\label{m39394*uid36}\item ${x}^{2}-x-20$
\newline
\newline
\end{enumerate}
        \label{m39394*uid37}\item Find the factors for the following trinomial expressions:
\label{m39394*id277980}\begin{enumerate}[noitemsep, label=\textbf{\alph*}. ] 
            \label{m39394*uid38}\item $2{x}^{2}+11x+5$\label{m39394*uid39}\item $3{x}^{2}+19x+6$\label{m39394*uid40}\item $6{x}^{2}+7x+2$\label{m39394*uid41}\item $12{x}^{2}+8x+1$\label{m39394*uid42}\item $8{x}^{2}+6x+1$
\newline
\newline
\end{enumerate}
        \label{m39394*uid43}\item Find the factors for the following trinomials:
\label{m39394*id278184}\begin{enumerate}[noitemsep, label=\textbf{\alph*}. ] 
            \label{m39394*uid44}\item $3{x}^{2}+17x-6$\label{m39394*uid45}\item $7{x}^{2}-6x-1$\label{m39394*uid46}\item $8{x}^{2}-6x+1$\label{m39394*uid47}\item $2{x}^{2}-5x-3$
\newline
\newline
\end{enumerate}
        \end{enumerate}
    \label{m39394*cid6}
\practiceinfo
\par 
 \par \begin{tabular}[h]{cccccc}
 (1.) 00iy&  (2.) 00iz&  (3.) 00j0&  (4.) 00j1& \end{tabular}
            \subsubsection{ Factorisation by Grouping}
            \nopagebreak
      \label{m39394*id278358}One other method of factorisation involves the use of common factors. We know that the factors of $3x+3$\hspace{1ex} are 3 and $\left(x+1\right)$. Similarly, the factors of $2{x}^{2}+2x$\hspace{1ex}are $2x$\hspace{1ex}and $\left(x+1\right)$. Therefore, if we have an expression:\par 
      \label{m39394*id278452}\nopagebreak\noindent{}
        \settowidth{\mymathboxwidth}{\begin{equation}
    2{x}^{2}+2x+3x+3\tag{8.30}
      \end{equation}
    }
    \typeout{Columnwidth = \the\columnwidth}\typeout{math as usual width = \the\mymathboxwidth}
    \ifthenelse{\lengthtest{\mymathboxwidth < \columnwidth}}{% if the math fits, do it again, for real
    \begin{equation}
    2{x}^{2}+2x+3x+3\tag{8.30}
      \end{equation}
    }{% else, if it doesn't fit
    \setlength{\mymathboxwidth}{\columnwidth}
      \addtolength{\mymathboxwidth}{-48pt}
    \par\vspace{12pt}\noindent\begin{minipage}{\columnwidth}
    \parbox[t]{\mymathboxwidth}{\large$
    2{x}^{2}+2x+3x+3$}\hfill
    \parbox[t]{48pt}{\raggedleft 
    (8.30)}
    \end{minipage}\vspace{12pt}\par
    }% end of conditional for this bit of math
    \typeout{math as usual width = \the\mymathboxwidth}
      \label{m39394*id278488}then we can factorise as:\par 
      \label{m39394*id278494}\nopagebreak\noindent{}
        \settowidth{\mymathboxwidth}{\begin{equation}
    2x\left(x+1\right)+3\left(x+1\right).\tag{8.31}
      \end{equation}
    }
    \typeout{Columnwidth = \the\columnwidth}\typeout{math as usual width = \the\mymathboxwidth}
    \ifthenelse{\lengthtest{\mymathboxwidth < \columnwidth}}{% if the math fits, do it again, for real
    \begin{equation}
    2x\left(x+1\right)+3\left(x+1\right).\tag{8.31}
      \end{equation}
    }{% else, if it doesn't fit
    \setlength{\mymathboxwidth}{\columnwidth}
      \addtolength{\mymathboxwidth}{-48pt}
    \par\vspace{12pt}\noindent\begin{minipage}{\columnwidth}
    \parbox[t]{\mymathboxwidth}{\large$
    2x\left(x+1\right)+3\left(x+1\right).$}\hfill
    \parbox[t]{48pt}{\raggedleft 
    (8.31)}
    \end{minipage}\vspace{12pt}\par
    }% end of conditional for this bit of math
    \typeout{math as usual width = \the\mymathboxwidth}
      \label{m39394*id278536}You can see that there is another common factor: $x+1$. Therefore, we can now write:\par 
      \label{m39394*id278556}\nopagebreak\noindent{}
        \settowidth{\mymathboxwidth}{\begin{equation}
    \left(x+1\right)\left(2x+3\right).\tag{8.32}
      \end{equation}
    }
    \typeout{Columnwidth = \the\columnwidth}\typeout{math as usual width = \the\mymathboxwidth}
    \ifthenelse{\lengthtest{\mymathboxwidth < \columnwidth}}{% if the math fits, do it again, for real
    \begin{equation}
    \left(x+1\right)\left(2x+3\right).\tag{8.32}
      \end{equation}
    }{% else, if it doesn't fit
    \setlength{\mymathboxwidth}{\columnwidth}
      \addtolength{\mymathboxwidth}{-48pt}
    \par\vspace{12pt}\noindent\begin{minipage}{\columnwidth}
    \parbox[t]{\mymathboxwidth}{\large$
    \left(x+1\right)\left(2x+3\right).$}\hfill
    \parbox[t]{48pt}{\raggedleft 
    (8.32)}
    \end{minipage}\vspace{12pt}\par
    }% end of conditional for this bit of math
    \typeout{math as usual width = \the\mymathboxwidth}
      \label{m39394*id278591}We get this by taking out the $x+1$ and seeing what is left over. We have a $+2x$\hspace{1ex}from the first term and a $+3$ from the second term. This is called \textsl{factorisation by grouping}.\par 
\label{m39394*secfhsst!!!underscore!!!id2835}\vspace{.5cm} 
      \noindent
      \hspace*{-30pt}\includegraphics[width=0.5in]{col11306.imgs/pspencil2.png}   \raisebox{25mm}{   
      \begin{mdframed}[linewidth=4, leftmargin=40, rightmargin=40]  
      \begin{exercise}
    \noindent\textbf{Exercise 8.9:  Factorisation by Grouping }
      \label{m39394*probfhsst!!!underscore!!!id2836}
      \label{m39394*id278656}Find the factors of $7x+14y+bx+2by$\hspace{1ex} by grouping \par 
      \vspace{5pt}
      \label{m39394*solfhsst!!!underscore!!!id2839}\noindent\textbf{Solution to Exercise } \label{m39394*listfhsst!!!underscore!!!id2839}\begin{enumerate}[noitemsep, label=\textbf{Step} \textbf{\arabic*}. ] 
            \leftskip=20pt\rightskip=\leftskip\item  
      \label{m39394*id278713}There are no factors that are common to all terms.\par 
      \item  
      \label{m39394*id278721}7 is a common factor of the first two terms and $b$ is a common factor of the second two terms.\par 
      \item  
      \label{m39394*id278739}\nopagebreak\noindent{}
        \settowidth{\mymathboxwidth}{\begin{equation}
    7x+14y+bx+2by=7\left(x+2y\right)+b\left(x+2y\right)\tag{8.33}
      \end{equation}
    }
    \typeout{Columnwidth = \the\columnwidth}\typeout{math as usual width = \the\mymathboxwidth}
    \ifthenelse{\lengthtest{\mymathboxwidth < \columnwidth}}{% if the math fits, do it again, for real
    \begin{equation}
    7x+14y+bx+2by=7\left(x+2y\right)+b\left(x+2y\right)\tag{8.33}
      \end{equation}
    }{% else, if it doesn't fit
    \setlength{\mymathboxwidth}{\columnwidth}
      \addtolength{\mymathboxwidth}{-48pt}
    \par\vspace{12pt}\noindent\begin{minipage}{\columnwidth}
    \parbox[t]{\mymathboxwidth}{\large$
    7x+14y+bx+2by=7\left(x+2y\right)+b\left(x+2y\right)$}\hfill
    \parbox[t]{48pt}{\raggedleft 
    (8.33)}
    \end{minipage}\vspace{12pt}\par
    }% end of conditional for this bit of math
    \typeout{math as usual width = \the\mymathboxwidth}
      \item  
      \label{m39394*id278811}$x+2y$\hspace{1ex}is a common factor.\par 
      \item  
      \label{m39394*id278835}\nopagebreak\noindent{}
        \settowidth{\mymathboxwidth}{\begin{equation}
    7\left(x+2y\right)+b\left(x+2y\right)=\left(x+2y\right)\left(7+b\right)\tag{8.34}
      \end{equation}
    }
    \typeout{Columnwidth = \the\columnwidth}\typeout{math as usual width = \the\mymathboxwidth}
    \ifthenelse{\lengthtest{\mymathboxwidth < \columnwidth}}{% if the math fits, do it again, for real
    \begin{equation}
    7\left(x+2y\right)+b\left(x+2y\right)=\left(x+2y\right)\left(7+b\right)\tag{8.34}
      \end{equation}
    }{% else, if it doesn't fit
    \setlength{\mymathboxwidth}{\columnwidth}
      \addtolength{\mymathboxwidth}{-48pt}
    \par\vspace{12pt}\noindent\begin{minipage}{\columnwidth}
    \parbox[t]{\mymathboxwidth}{\large$
    7\left(x+2y\right)+b\left(x+2y\right)=\left(x+2y\right)\left(7+b\right)$}\hfill
    \parbox[t]{48pt}{\raggedleft 
    (8.34)}
    \end{minipage}\vspace{12pt}\par
    }% end of conditional for this bit of math
    \typeout{math as usual width = \the\mymathboxwidth}
      \item  
      \label{m39394*id278906}The factors of $7x+14y+bx+2by$\hspace{1ex}are $\left(7+b\right)$ and $\left(x+2y\right)$.
 \par 
      \end{enumerate}
    \end{exercise}
    \end{mdframed}
    }
    \noindent
\label{m39394*eip-280}
    \setcounter{subfigure}{0}
	\begin{figure}[H] % horizontal\label{m39394*factorisingtrinomial}
    \textnormal{Khan Academy video on factorising a trinomial by grouping.}\vspace{.1in} \nopagebreak
  \label{m39394*yt-media32}\label{m39394*yt-video32}
            \raisebox{-5 pt}{ \includegraphics[width=0.5cm]{col11306.imgs/summary_www.png}} { (Video:  MG10065 )}
      \vspace{2pt}
    \vspace{.1in}
 \end{figure}       \par \label{m39394*secfhsst!!!underscore!!!id2920}
            \subsubsection{  Factorisation by Grouping }
            \nopagebreak
      \label{m39394*id279000}\begin{enumerate}[noitemsep, label=\textbf{\arabic*}. ] 
            \label{m39394*uid48}\item Factorise by grouping: $6x+a+2ax+3$
\newline
        \label{m39394*uid49}\item Factorise by grouping: ${x}^{2}-6x+5x-30$\newline
        \label{m39394*uid50}\item Factorise by grouping: $5x+10y-ax-2ay$\newline
        \label{m39394*uid51}\item Factorise by grouping: ${a}^{2}-2a-ax+2x$\newline
        \label{m39394*uid52}\item Factorise by grouping: $5xy-3y+10x-6$\newline
        \end{enumerate}
  \label{m39394**end}
\practiceinfo
\par 
 \par \begin{tabular}[h]{cccccc}
 (1.) 00j2&  (2.) 00j3&  (3.) 00j4&  (4.) 00j5&  (5.) 00j6& \end{tabular}
         \subsection{ Fractions}
    \nopagebreak
            \label{m39392} $ \hspace{-5pt}\begin{array}{cccccccccccc}   \includegraphics[width=0.75cm]{col11306.imgs/summary_fullmarks.png} &   \end{array} $ \hspace{2 pt}\raisebox{-5 pt}{} {(subsection shortcode: MG10066 )} \par 
    \label{m39392*cid7}
            \subsubsection{ Simplification of Fractions}
            \nopagebreak
      \label{m39392*id279238}In some cases of simplifying an algebraic expression, the expression will be a fraction. For example,\par 
      \label{m39392*id279242}\nopagebreak\noindent{}
        \settowidth{\mymathboxwidth}{\begin{equation}
    \frac{{x}^{2}+3x}{x+3}\tag{8.35}
      \end{equation}
    }
    \typeout{Columnwidth = \the\columnwidth}\typeout{math as usual width = \the\mymathboxwidth}
    \ifthenelse{\lengthtest{\mymathboxwidth < \columnwidth}}{% if the math fits, do it again, for real
    \begin{equation}
    \frac{{x}^{2}+3x}{x+3}\tag{8.35}
      \end{equation}
    }{% else, if it doesn't fit
    \setlength{\mymathboxwidth}{\columnwidth}
      \addtolength{\mymathboxwidth}{-48pt}
    \par\vspace{12pt}\noindent\begin{minipage}{\columnwidth}
    \parbox[t]{\mymathboxwidth}{\large$
    \frac{{x}^{2}+3x}{x+3}$}\hfill
    \parbox[t]{48pt}{\raggedleft 
    (8.35)}
    \end{minipage}\vspace{12pt}\par
    }% end of conditional for this bit of math
    \typeout{math as usual width = \the\mymathboxwidth}
      \label{m39392*id279276}has a quadratic in the numerator and a binomial in the denominator. You can apply the different factorisation methods to simplify the expression.\par 
      \label{m39392*id279282}\nopagebreak\noindent{}\settowidth{\mymathboxwidth}{\begin{equation}
    \begin{array}{cccc}& \phantom{\rule{4pt}{0ex}}& \frac{{x}^{2}+3x}{x+3}\hfill & \\ & =& \frac{x\left(x+3\right)}{x+3}\hfill & \\ & =& x\hfill & \mathrm{provided}\phantom{\rule{2pt}{0ex}}x\ne -3\hfill \end{array}\tag{8.36}
      \end{equation}
    }
    \typeout{Columnwidth = \the\columnwidth}\typeout{math as usual width = \the\mymathboxwidth}
    \ifthenelse{\lengthtest{\mymathboxwidth < \columnwidth}}{% if the math fits, do it again, for real
    \begin{equation}
    \begin{array}{cccc}& \phantom{\rule{4pt}{0ex}}& \frac{{x}^{2}+3x}{x+3}\hfill & \\ & =& \frac{x\left(x+3\right)}{x+3}\hfill & \\ & =& x\hfill & \mathrm{provided}\phantom{\rule{2pt}{0ex}}x\ne -3\hfill \end{array}\tag{8.36}
      \end{equation}
    }{% else, if it doesn't fit
    \setlength{\mymathboxwidth}{\columnwidth}
      \addtolength{\mymathboxwidth}{-48pt}
    \par\vspace{12pt}\noindent\begin{minipage}{\columnwidth}
    \parbox[t]{\mymathboxwidth}{\large$
    \phantom{\rule{4pt}{0ex}}\frac{{x}^{2}+3x}{x+3}=\frac{x\left(x+3\right)}{x+3}=x\mathrm{provided}\phantom{\rule{2pt}{0ex}}x\ne -3$}\hfill
    \parbox[t]{48pt}{\raggedleft 
    (8.36)}
    \end{minipage}\vspace{12pt}\par
    }% end of conditional for this bit of math
    \typeout{math as usual width = \the\mymathboxwidth}
      \label{m39392*id279389}If $x$ were 3 then the denominator, $x-3$, would be 0 and the fraction undefined.\par 
\label{m39392*secfhsst!!!underscore!!!id3026}\vspace{.5cm} 
      \noindent
      \hspace*{-30pt}\includegraphics[width=0.5in]{col11306.imgs/pspencil2.png}   \raisebox{25mm}{   
      \begin{mdframed}[linewidth=4, leftmargin=40, rightmargin=40]  
      \begin{exercise}
    \noindent\textbf{Exercise 8.10:  Simplification of Fractions }
      \label{m39392*probfhsst!!!underscore!!!id3027}
      \label{m39392*id279431}Simplify: $\frac{2x-b+x-ab}{a{x}^{2}-abx}$ \par 
      \vspace{5pt}
      \label{m39392*solfhsst!!!underscore!!!id3030}\noindent\textbf{Solution to Exercise } \label{m39392*listfhsst!!!underscore!!!id3030}\begin{enumerate}[noitemsep, label=\textbf{Step} \textbf{\arabic*}. ] 
            \leftskip=20pt\rightskip=\leftskip\item  
      \label{m39392*id279503}Use \textsl{grouping} for numerator and \textsl{common factor} for denominator in this example.\par 
      \label{m39392*id279518}\nopagebreak\noindent{}
        \settowidth{\mymathboxwidth}{\begin{equation}
    \begin{array}{ccc}& =& \frac{\left(ax-ab\right)+\left(x-b\right)}{a{x}^{2}-abx}\hfill \\ & =& \frac{a\left(x-b\right)+\left(x-b\right)}{ax\left(x-b\right)}\hfill \\ & =& \frac{\left(x-b\right)\left(a+1\right)}{ax\left(x-b\right)}\hfill \end{array}\tag{8.37}
      \end{equation}
    }
    \typeout{Columnwidth = \the\columnwidth}\typeout{math as usual width = \the\mymathboxwidth}
    \ifthenelse{\lengthtest{\mymathboxwidth < \columnwidth}}{% if the math fits, do it again, for real
    \begin{equation}
    \begin{array}{ccc}& =& \frac{\left(ax-ab\right)+\left(x-b\right)}{a{x}^{2}-abx}\hfill \\ & =& \frac{a\left(x-b\right)+\left(x-b\right)}{ax\left(x-b\right)}\hfill \\ & =& \frac{\left(x-b\right)\left(a+1\right)}{ax\left(x-b\right)}\hfill \end{array}\tag{8.37}
      \end{equation}
    }{% else, if it doesn't fit
    \setlength{\mymathboxwidth}{\columnwidth}
      \addtolength{\mymathboxwidth}{-48pt}
    \par\vspace{12pt}\noindent\begin{minipage}{\columnwidth}
    \parbox[t]{\mymathboxwidth}{\large$
    =\frac{\left(ax-ab\right)+\left(x-b\right)}{a{x}^{2}-abx}=\frac{a\left(x-b\right)+\left(x-b\right)}{ax\left(x-b\right)}=\frac{\left(x-b\right)\left(a+1\right)}{ax\left(x-b\right)}$}\hfill
    \parbox[t]{48pt}{\raggedleft 
    (8.37)}
    \end{minipage}\vspace{12pt}\par
    }% end of conditional for this bit of math
    \typeout{math as usual width = \the\mymathboxwidth}
      \item  
      \label{m39392*id279695}The simplified answer is:\par 
      \label{m39392*id279699}\nopagebreak\noindent{}
        \settowidth{\mymathboxwidth}{\begin{equation}
    \begin{array}{ccc}& =& \frac{a+1}{ax}\hfill \end{array}\tag{8.38}
      \end{equation}
    }
    \typeout{Columnwidth = \the\columnwidth}\typeout{math as usual width = \the\mymathboxwidth}
    \ifthenelse{\lengthtest{\mymathboxwidth < \columnwidth}}{% if the math fits, do it again, for real
    \begin{equation}
    \begin{array}{ccc}& =& \frac{a+1}{ax}\hfill \end{array}\tag{8.38}
      \end{equation}
    }{% else, if it doesn't fit
    \setlength{\mymathboxwidth}{\columnwidth}
      \addtolength{\mymathboxwidth}{-48pt}
    \par\vspace{12pt}\noindent\begin{minipage}{\columnwidth}
    \parbox[t]{\mymathboxwidth}{\large$
    =\frac{a+1}{ax}$}\hfill
    \parbox[t]{48pt}{\raggedleft 
    (8.38)}
    \end{minipage}\vspace{12pt}\par
    }% end of conditional for this bit of math
    \typeout{math as usual width = \the\mymathboxwidth}
      \end{enumerate}
    \end{exercise}
    \end{mdframed}
    }
    \noindent
\label{m39392*secfhsst!!!underscore!!!id3174}\vspace{.5cm} 
      \noindent
      \hspace*{-30pt}\includegraphics[width=0.5in]{col11306.imgs/pspencil2.png}   \raisebox{25mm}{   
      \begin{mdframed}[linewidth=4, leftmargin=40, rightmargin=40]  
      \begin{exercise}
    \noindent\textbf{Exercise 8.11:  Simplification of Fractions }
      \label{m39392*probfhsst!!!underscore!!!id3175}
      \label{m39392*id279776}Simplify:$\frac{{x}^{2}-x-2}{{x}^{2}-4}÷\frac{{x}^{2}+x}{{x}^{2}+2x}$ \par 
      \vspace{5pt}
      \label{m39392*solfhsst!!!underscore!!!id3178}\noindent\textbf{Solution to Exercise } \label{m39392*listfhsst!!!underscore!!!id3178}\begin{enumerate}[noitemsep, label=\textbf{Step} \textbf{\arabic*}. ] 
            \leftskip=20pt\rightskip=\leftskip\item  
      \label{m39392*id279872}\nopagebreak\noindent{}
        \settowidth{\mymathboxwidth}{\begin{equation}
    \begin{array}{ccc}& =& \frac{\left(x+1\right)\left(x-2\right)}{\left(x+2\right)\left(x-2\right)}÷\frac{x\left(x+1\right)}{x\left(x+2\right)}\hfill \end{array}\tag{8.39}
      \end{equation}
    }
    \typeout{Columnwidth = \the\columnwidth}\typeout{math as usual width = \the\mymathboxwidth}
    \ifthenelse{\lengthtest{\mymathboxwidth < \columnwidth}}{% if the math fits, do it again, for real
    \begin{equation}
    \begin{array}{ccc}& =& \frac{\left(x+1\right)\left(x-2\right)}{\left(x+2\right)\left(x-2\right)}÷\frac{x\left(x+1\right)}{x\left(x+2\right)}\hfill \end{array}\tag{8.39}
      \end{equation}
    }{% else, if it doesn't fit
    \setlength{\mymathboxwidth}{\columnwidth}
      \addtolength{\mymathboxwidth}{-48pt}
    \par\vspace{12pt}\noindent\begin{minipage}{\columnwidth}
    \parbox[t]{\mymathboxwidth}{\large$
    =\frac{\left(x+1\right)\left(x-2\right)}{\left(x+2\right)\left(x-2\right)}÷\frac{x\left(x+1\right)}{x\left(x+2\right)}$}\hfill
    \parbox[t]{48pt}{\raggedleft 
    (8.39)}
    \end{minipage}\vspace{12pt}\par
    }% end of conditional for this bit of math
    \typeout{math as usual width = \the\mymathboxwidth}
      \item  
      \label{m39392*id279977}\nopagebreak\noindent{}
        \settowidth{\mymathboxwidth}{\begin{equation}
    \begin{array}{ccc}& =& \frac{\left(x+1\right)\left(x-2\right)}{\left(x+2\right)\left(x-2\right)}\ensuremath{\times}\frac{x\left(x+2\right)}{x\left(x+1\right)}\hfill \end{array}\tag{8.40}
      \end{equation}
    }
    \typeout{Columnwidth = \the\columnwidth}\typeout{math as usual width = \the\mymathboxwidth}
    \ifthenelse{\lengthtest{\mymathboxwidth < \columnwidth}}{% if the math fits, do it again, for real
    \begin{equation}
    \begin{array}{ccc}& =& \frac{\left(x+1\right)\left(x-2\right)}{\left(x+2\right)\left(x-2\right)}\ensuremath{\times}\frac{x\left(x+2\right)}{x\left(x+1\right)}\hfill \end{array}\tag{8.40}
      \end{equation}
    }{% else, if it doesn't fit
    \setlength{\mymathboxwidth}{\columnwidth}
      \addtolength{\mymathboxwidth}{-48pt}
    \par\vspace{12pt}\noindent\begin{minipage}{\columnwidth}
    \parbox[t]{\mymathboxwidth}{\large$
    =\frac{\left(x+1\right)\left(x-2\right)}{\left(x+2\right)\left(x-2\right)}\ensuremath{\times}\frac{x\left(x+2\right)}{x\left(x+1\right)}$}\hfill
    \parbox[t]{48pt}{\raggedleft 
    (8.40)}
    \end{minipage}\vspace{12pt}\par
    }% end of conditional for this bit of math
    \typeout{math as usual width = \the\mymathboxwidth}
      \item  
      \label{m39392*id280081}The simplified answer is\par 
      \label{m39392*id280085}\nopagebreak\noindent{}
        \settowidth{\mymathboxwidth}{\begin{equation}
    \begin{array}{ccc}& =& 1\hfill \end{array}\tag{8.41}
      \end{equation}
    }
    \typeout{Columnwidth = \the\columnwidth}\typeout{math as usual width = \the\mymathboxwidth}
    \ifthenelse{\lengthtest{\mymathboxwidth < \columnwidth}}{% if the math fits, do it again, for real
    \begin{equation}
    \begin{array}{ccc}& =& 1\hfill \end{array}\tag{8.41}
      \end{equation}
    }{% else, if it doesn't fit
    \setlength{\mymathboxwidth}{\columnwidth}
      \addtolength{\mymathboxwidth}{-48pt}
    \par\vspace{12pt}\noindent\begin{minipage}{\columnwidth}
    \parbox[t]{\mymathboxwidth}{\large$
    =1$}\hfill
    \parbox[t]{48pt}{\raggedleft 
    (8.41)}
    \end{minipage}\vspace{12pt}\par
    }% end of conditional for this bit of math
    \typeout{math as usual width = \the\mymathboxwidth}
      \end{enumerate}
    \end{exercise}
    \end{mdframed}
    }
    \noindent
\label{m39392*secfhsst!!!underscore!!!id3323}
            \subsubsection{  Simplification of Fractions }
            \nopagebreak
      \label{m39392*id280133}\begin{enumerate}[noitemsep, label=\textbf{\arabic*}. ] 
            \label{m39392*uid53}\item Simplify:
    % \textbf{m39392*id280148}\par
          \begin{table}[H]
    % \begin{table}[H]
    % \\ 'id2896760' '1'
        \begin{center}
      \label{m39392*id280148}
    \noindent
    \tabletail{%
        \hline
        \multicolumn{2}{|p{\mytableboxwidth}|}{\raggedleft \small \sl continued on next page}\\
        \hline
      }
      \tablelasttail{}
      \begin{xtabular}[t]{|l|l|}\hline
        (a) $\frac{3a}{15}$\hspace{1ex} &
        (b) $\frac{2a+10}{4}$\hspace{1ex}% make-rowspan-placeholders
     \tabularnewline\cline{1-1}\cline{2-2}
      %--------------------------------------------------------------------
        (c) $\frac{5a+20}{a+4}$\hspace{1ex} &
        (d) $\frac{{a}^{2}-4a}{a-4}$\hspace{1ex}% make-rowspan-placeholders
     \tabularnewline\cline{1-1}\cline{2-2}
      %--------------------------------------------------------------------
        (e) $\frac{3{a}^{2}-9a}{2a-6}$\hspace{1ex} &
        (f) $\frac{9a+27}{9a+18}$\hspace{1ex}% make-rowspan-placeholders
     \tabularnewline\cline{1-1}\cline{2-2}
      %--------------------------------------------------------------------
        (g) $\frac{6ab+2a}{2b}$\hspace{1ex} &
        (h) $\frac{16{x}^{2}y-8xy}{12x-6}$\hspace{1ex}% make-rowspan-placeholders
     \tabularnewline\cline{1-1}\cline{2-2}
      %--------------------------------------------------------------------
        (i) $\frac{4xyp-8xp}{12xy}$\hspace{1ex} &
        (j) $\frac{3a+9}{14}÷\frac{7a+21}{a+3}$\hspace{1ex}% make-rowspan-placeholders
     \tabularnewline\cline{1-1}\cline{2-2}
      %--------------------------------------------------------------------
        (k) $\frac{{a}^{2}-5a}{2a+10}÷\frac{3a+15}{4a}$\hspace{1ex} &
        (l) $\frac{3xp+4p}{8p}÷\frac{12{p}^{2}}{3x+4}$\hspace{1ex}% make-rowspan-placeholders
     \tabularnewline\cline{1-1}\cline{2-2}
      %--------------------------------------------------------------------
        (m) $\frac{16}{2xp+4x}÷\frac{6{x}^{2}+8x}{12}$\hspace{1ex} &
        (n) $\frac{24a-8}{12}÷\frac{9a-3}{6}$\hspace{1ex}% make-rowspan-placeholders
     \tabularnewline\cline{1-1}\cline{2-2}
      %--------------------------------------------------------------------
        (o) $\frac{{a}^{2}+2a}{5}÷\frac{2a+4}{20}$\hspace{1ex} &
        (p) $\frac{{p}^{2}+pq}{7p}÷\frac{8p+8q}{21q}$\hspace{1ex}% make-rowspan-placeholders
     \tabularnewline\cline{1-1}\cline{2-2}
      %--------------------------------------------------------------------
        (q) $\frac{5ab-15b}{4a-12}÷\frac{6{b}^{2}}{a+b}$\hspace{1ex} &
        (r) $\frac{{f}^{2}a-f{a}^{2}}{f-a}$% make-rowspan-placeholders
     \tabularnewline\cline{1-1}\cline{2-2}
      %--------------------------------------------------------------------
    \end{xtabular}
      \end{center}
    \begin{center}{\small\bfseries Table 8.11}\end{center}
    \begin{caption}{\small\bfseries Table 8.11}\end{caption}
\end{table}
    \par
          \label{m39392*uid54}\item Simplify: $\frac{{x}^{2}-1}{3}\ensuremath{\times}\frac{1}{x-1}-\frac{1}{2}$
\newline
\newline
\end{enumerate}
    \label{m39392*eip-770}
\practiceinfo
\par 
 \par \begin{tabular}[h]{cccccc}
 (1.) 00j7&  (2.) 00j8& \end{tabular}
            \subsubsection{ Adding and subtracting fractions}
            \nopagebreak
            \label{m39392*eip-107}Using the concepts learnt in simplification of fractions, we can now add and subtract simple fractions. To add or subtract fractions we note that we can only add or subtract fractions that have the same denominator. So we must first make all the denominators the same and then perform the addition or subtraction. This is called finding the lowest common denominator or multiple. 
\par \label{m39392*eip-997}For example, if you wanted to add: $\frac{1}{2}$ and $\frac{3}{5}$ we would note that the lowest common denominator is 10. So we must multiply the first fraction by 5 and the second fraction by 2 to get both of these with the same denominator. Doing so gives: $\frac{5}{10}$ and $\frac{6}{10}$. Now we can add the fractions. Doing so, we get $\frac{11}{10}$.\par \label{m39392*eip-337}\vspace{.5cm} 
      \noindent
      \hspace*{-30pt}\includegraphics[width=0.5in]{col11306.imgs/pspencil2.png}   \raisebox{25mm}{   
      \begin{mdframed}[linewidth=4, leftmargin=40, rightmargin=40]  
      \begin{exercise}
    \noindent\textbf{Exercise 8.12}\label{m39392*probfhsst!!!underscore!!!id317245}
      \label{m39392*id2793436776}Simplify the following expression: $\frac{x-2}{{x}^{2}-4}+\frac{{x}^{2}}{x-2}-\frac{{x}^{3}+x-4}{{x}^{2}-4}$ \par 
      \vspace{5pt}
      \label{m39392*solfhsst!!!underscore!!!id3132478}\noindent\textbf{Solution to Exercise } \label{m39392*listfhsst!!!underscore!!!id343178}\begin{enumerate}[noitemsep, label=\textbf{Step} \textbf{\arabic*}. ] 
            \leftskip=20pt\rightskip=\leftskip\item  
      \label{m39392*id279353872}\nopagebreak\noindent{}\settowidth{\mymathboxwidth}{\begin{equation}
    \frac{x-2}{\left(x+2\right)\left(x-2\right)}+\frac{{x}^{2}}{x-2}-\frac{{x}^{3}+x-4}{\left(x+2\right)\left(x-2\right)}\tag{8.42}
      \end{equation}
    }
    \typeout{Columnwidth = \the\columnwidth}\typeout{math as usual width = \the\mymathboxwidth}
    \ifthenelse{\lengthtest{\mymathboxwidth < \columnwidth}}{% if the math fits, do it again, for real
    \begin{equation}
    \frac{x-2}{\left(x+2\right)\left(x-2\right)}+\frac{{x}^{2}}{x-2}-\frac{{x}^{3}+x-4}{\left(x+2\right)\left(x-2\right)}\tag{8.42}
      \end{equation}
    }{% else, if it doesn't fit
    \setlength{\mymathboxwidth}{\columnwidth}
      \addtolength{\mymathboxwidth}{-48pt}
    \par\vspace{12pt}\noindent\begin{minipage}{\columnwidth}
    \parbox[t]{\mymathboxwidth}{\large$
    \frac{x-2}{\left(x+2\right)\left(x-2\right)}+\frac{{x}^{2}}{x-2}-\frac{{x}^{3}+x-4}{\left(x+2\right)\left(x-2\right)}$}\hfill
    \parbox[t]{48pt}{\raggedleft 
    (8.42)}
    \end{minipage}\vspace{12pt}\par
    }% end of conditional for this bit of math
    \typeout{math as usual width = \the\mymathboxwidth}
      \item  \label{m39392*id79242}We make all the denominators the same so that we can add or subtract the fractions. The lowest common denominator is $\left(x-2\right)\left(x+2\right)$.\par 
      \label{m39392*id27943653977}\nopagebreak\noindent{}\settowidth{\mymathboxwidth}{\begin{equation}
    \frac{x-2}{\left(x+2\right)\left(x-2\right)}+\frac{\left({x}^{2}\right)
\left(x+2\right)}{\left(x+2\right)\left(x-2\right)}-\frac{{x}^{3}+x-4}{\left(x+2\right)\left(x-2\right)}\tag{8.43}
      \end{equation}
    }
    \typeout{Columnwidth = \the\columnwidth}\typeout{math as usual width = \the\mymathboxwidth}
    \ifthenelse{\lengthtest{\mymathboxwidth < \columnwidth}}{% if the math fits, do it again, for real
    \begin{equation}
    \frac{x-2}{\left(x+2\right)\left(x-2\right)}+\frac{\left({x}^{2}\right)
\left(x+2\right)}{\left(x+2\right)\left(x-2\right)}-\frac{{x}^{3}+x-4}{\left(x+2\right)\left(x-2\right)}\tag{8.43}
      \end{equation}
    }{% else, if it doesn't fit
    \setlength{\mymathboxwidth}{\columnwidth}
      \addtolength{\mymathboxwidth}{-48pt}
    \par\vspace{12pt}\noindent\begin{minipage}{\columnwidth}
    \parbox[t]{\mymathboxwidth}{\large$
    \frac{x-2}{\left(x+2\right)\left(x-2\right)}+\frac{\left({x}^{2}\right)
\left(x+2\right)}{\left(x+2\right)\left(x-2\right)}-\frac{{x}^{3}+x-4}{\left(x+2\right)\left(x-2\right)}$}\hfill
    \parbox[t]{48pt}{\raggedleft 
    (8.43)}
    \end{minipage}\vspace{12pt}\par
    }% end of conditional for this bit of math
    \typeout{math as usual width = \the\mymathboxwidth}
      \item  \label{m39392*id639247}Since the fractions all have the same denominator we can write them all as one fraction with the appropriate operator\par 
     \label{m39392*id653977}\nopagebreak\noindent{}\settowidth{\mymathboxwidth}{\begin{equation}
    \frac{x-2+\left({x}^{2}\right)
\left(x+2\right)-{x}^{3}+x-4}{\left(x+2\right)\left(x-2\right)}\tag{8.44}
      \end{equation}
    }
    \typeout{Columnwidth = \the\columnwidth}\typeout{math as usual width = \the\mymathboxwidth}
    \ifthenelse{\lengthtest{\mymathboxwidth < \columnwidth}}{% if the math fits, do it again, for real
    \begin{equation}
    \frac{x-2+\left({x}^{2}\right)
\left(x+2\right)-{x}^{3}+x-4}{\left(x+2\right)\left(x-2\right)}\tag{8.44}
      \end{equation}
    }{% else, if it doesn't fit
    \setlength{\mymathboxwidth}{\columnwidth}
      \addtolength{\mymathboxwidth}{-48pt}
    \par\vspace{12pt}\noindent\begin{minipage}{\columnwidth}
    \parbox[t]{\mymathboxwidth}{\large$
    \frac{x-2+\left({x}^{2}\right)
\left(x+2\right)-{x}^{3}+x-4}{\left(x+2\right)\left(x-2\right)}$}\hfill
    \parbox[t]{48pt}{\raggedleft 
    (8.44)}
    \end{minipage}\vspace{12pt}\par
    }% end of conditional for this bit of math
    \typeout{math as usual width = \the\mymathboxwidth}
      \item \label{m39392*id9657}\nopagebreak\noindent{}\settowidth{\mymathboxwidth}{\begin{equation}
    \frac{2{x}^{2}
+2x-6}{\left(x+2\right)\left(x-2\right)}\tag{8.45}
      \end{equation}
    }
    \typeout{Columnwidth = \the\columnwidth}\typeout{math as usual width = \the\mymathboxwidth}
    \ifthenelse{\lengthtest{\mymathboxwidth < \columnwidth}}{% if the math fits, do it again, for real
    \begin{equation}
    \frac{2{x}^{2}
+2x-6}{\left(x+2\right)\left(x-2\right)}\tag{8.45}
      \end{equation}
    }{% else, if it doesn't fit
    \setlength{\mymathboxwidth}{\columnwidth}
      \addtolength{\mymathboxwidth}{-48pt}
    \par\vspace{12pt}\noindent\begin{minipage}{\columnwidth}
    \parbox[t]{\mymathboxwidth}{\large$
    \frac{2{x}^{2}
+2x-6}{\left(x+2\right)\left(x-2\right)}$}\hfill
    \parbox[t]{48pt}{\raggedleft 
    (8.45)}
    \end{minipage}\vspace{12pt}\par
    }% end of conditional for this bit of math
    \typeout{math as usual width = \the\mymathboxwidth}
\item 
\label{m39392*id65677}\nopagebreak\noindent{}\settowidth{\mymathboxwidth}{\begin{equation}
    \frac{2\left({x}^{2}
+x-3\right)}{\left(x+2\right)\left(x-2\right)}\tag{8.46}
      \end{equation}
    }
    \typeout{Columnwidth = \the\columnwidth}\typeout{math as usual width = \the\mymathboxwidth}
    \ifthenelse{\lengthtest{\mymathboxwidth < \columnwidth}}{% if the math fits, do it again, for real
    \begin{equation}
    \frac{2\left({x}^{2}
+x-3\right)}{\left(x+2\right)\left(x-2\right)}\tag{8.46}
      \end{equation}
    }{% else, if it doesn't fit
    \setlength{\mymathboxwidth}{\columnwidth}
      \addtolength{\mymathboxwidth}{-48pt}
    \par\vspace{12pt}\noindent\begin{minipage}{\columnwidth}
    \parbox[t]{\mymathboxwidth}{\large$
    \frac{2\left({x}^{2}
+x-3\right)}{\left(x+2\right)\left(x-2\right)}$}\hfill
    \parbox[t]{48pt}{\raggedleft 
    (8.46)}
    \end{minipage}\vspace{12pt}\par
    }% end of conditional for this bit of math
    \typeout{math as usual width = \the\mymathboxwidth}
\end{enumerate}
    \end{exercise}
    \end{mdframed}
    }
    \noindent
  \label{m39392*eip-755}
            \subsubsection{ Two interesting mathematical proofs}
            \nopagebreak
            \label{m39392*id76534}We can use the concepts learnt in this chapter to demonstrate two interesting mathematical proofs. The first proof states that ${n}^{2}+n$ is even for all $n\in {Z}$. The second proof states that ${n}^{3}-n$ is divisible by 6 for all $n\in {Z}$. Before we demonstrate that these two laws are true, we first need to note some other mathematical rules.
\par 
\label{m39392*id736}If we multiply an even number by an odd number, we get an even number. Similarly if we multiply an odd number by an even number we get an even number. Also, an even number multiplied by an even number is even and an odd number multiplied by an odd number is odd. This result is shown in the following table:
\par 
    % \textbf{m39392*eip-556}\par
          \begin{table}[H]
    % \begin{table}[H]
    % \\ '' '0'
        \begin{center}
      \label{m39392*eip-556}
    \noindent
    \tabletail{%
        \hline
        \multicolumn{3}{|p{\mytableboxwidth}|}{\raggedleft \small \sl continued on next page}\\
        \hline
      }
      \tablelasttail{}
      \begin{xtabular}[t]{|l|l|l|}\hline
         &
        \textbf{Odd number} &
        \textbf{Even number}% make-rowspan-placeholders
     \tabularnewline\cline{1-1}\cline{2-2}\cline{3-3}
      %--------------------------------------------------------------------
        \textbf{Odd number} &
        Odd &
        Even% make-rowspan-placeholders
     \tabularnewline\cline{1-1}\cline{2-2}\cline{3-3}
      %--------------------------------------------------------------------
        \textbf{Even number} &
        Even &
        Even% make-rowspan-placeholders
     \tabularnewline\cline{1-1}\cline{2-2}\cline{3-3}
      %--------------------------------------------------------------------
    \end{xtabular}
      \end{center}
    \begin{center}{\small\bfseries Table 8.12}\end{center}
    \begin{caption}{\small\bfseries Table 8.12}\end{caption}
\end{table}
    \par
  \label{m39392*id697}If we take three consecutive numbers and multiply them together, the resulting number is always divisible by three. This should be obvious since if we have any three consecutive numbers, one of them will be divisible by 3.
\par 
\label{m39392*id67354}Now we are ready to demonstrate that ${n}^{2}+n$ is even for all $n\in {Z}$. If we factorise this expression we get:
$n\left(n+1\right)$. If $n$ is even, than $n+1$ is odd. If $n$ is odd, than $n+1$ is even. Since we know that if we multiply an even number with an odd number or an odd number with an even number, we get an even number, we have demonstrated that ${n}^{2}+n$ is always even. Try this for a few values of $n$ and you should find that this is true.
\par 
\label{m39392*id67744}To demonstrate that ${n}^{3}-n$ is divisible by 6 for all $n\in {Z}$, we first note that the factors of 6 are 3 and 2. So if we show that ${n}^{3}-n$ is divisible by both 3 and 2, then we have shown that it is also divisible by 6! If we factorise this expression we get:
$n\left(n+1\right)\left(n-1\right)$. Now we note that we are multiplying three consecutive numbers together (we are taking $n$ and then adding 1 or subtracting 1. This gives us the two numbers on either side of $n$.) For example, if $n=4$, then $n+1=5$ and $n-1=3$. But we know that when we multiply three consecutive numbers together, the resulting number is always divisible by 3. So we have demonstrated that ${n}^{3}-n$ is always divisible by 3. To demonstrate that it is also divisible by 2, we can also show that it is even. We have shown that ${n}^{2}+n$ is always even. So now we recall what we said about multiplying even and odd numbers. Since one number is always even and the other can be either even or odd, the result of multiplying these numbers together is always even. And so we have demonstrated that ${n}^{3}-n$ is divisible by 6 for all $n\in {Z}$.
\par 
\label{m39392*eip-735}
            \subsubsection{ Summary}
            \nopagebreak
            \label{m39392*uid0812}\begin{itemize}[noitemsep]
            \item A binomial is a mathematical expression with two terms. The product of two identical binomials is known as the square of the binomial. The difference of two squares is when we multiply
                $\left(ax+b\right)\left(ax-b\right)$\item Factorising is the opposite of expanding the brackets. You can use common factors or the difference of two squares to help you factorise expressions.\item The distributive law ($\left(A+B\right)\left(C+D+E\right)=A\left(C+D+E\right)+B\left(C+D+E\right)$) helps us to multiply a binomial and a trinomial.\item The sum of cubes is: $\left(x+y\right)\left({x}^{2}-xy+{y}^{2}\right)={x}^{3}+{y}^{3}$ and the difference of cubes is: ${x}^{3}-{y}^{3}=\left(x-y\right)\left({x}^{2}+xy+{y}^{2}\right)$\item To factorise a quadratic we find the two binomials that were multiplied together to give the quadratic.\item We can also factorise a quadratic by grouping. This is where we find a common factor in the quadratic and take it out and then see what is left over.\item We can simplify fractions by using the methods we have learnt to factorise expressions.\item Fractions can be added or subtracted. To do this the denominators of each fraction must be the same.\end{itemize}
        \label{m39392*cid8}
            \subsubsection{ End of Chapter Exercises}
            \nopagebreak
      \label{m39392*id281011}\begin{enumerate}[noitemsep, label=\textbf{\arabic*}. ] 
            \label{m39392*uid55}\item Factorise:
\label{m39392*id281026}\begin{enumerate}[noitemsep, label=\textbf{\alph*}. ] 
            \item ${a}^{2}-9$\item ${m}^{2}-36$\item $9{b}^{2}-81$\item $16{b}^{6}-25{a}^{2}$\item ${m}^{2}-\left(1/9\right)$\item $5-5{a}^{2}{b}^{6}$\item $16b{a}^{4}-81b$\item ${a}^{2}-10a+25$\item $16{b}^{2}+56b+49$\item $2{a}^{2}-12ab+18{b}^{2}$\item $-4{b}^{2}-144{b}^{8}+48{b}^{5}$\end{enumerate}
                \label{m39392*id632}\item Factorise completely: \label{m39392*id6423}\begin{enumerate}[noitemsep, label=\textbf{\alph*}. ] 
            \item $\left(16-{x}^{4}\right)$\item ${7x}^{2}-14x+7xy-14y$
\item ${y}^{2}-7y-30$
\item $1-x-{x}^{2}+{x}^{3}$
\item $-3\left(1-{p}^{2}\right)+p+1$\end{enumerate}
\item Simplify the following:
\label{m39392*eip-id1166762435067}\begin{enumerate}[noitemsep, label=\textbf{\alph*}. ] 
            \item ${\left(a-2\right)}^{2}-a\left(a+4\right)$\item $\left(5a-4b\right)\left(25{a}^{2}+20\mathrm{ab}+16{b}^{2}\right)$\item $\left(2m-3\right)\left(4{m}^{2}+9\right)\left(2m+3\right)$\item $\left(a+2b-c\right)\left(a+2b+c\right)$\end{enumerate}
\item Simplify the following:
\label{m39392*eip-id1153}\begin{enumerate}[noitemsep, label=\textbf{\alph*}. ] 
            \item $\frac{{p}^{2}-{q}^{2}}{p}÷\frac{p+q}{{p}^{2}-\mathrm{pq}}$\item $\frac{2}{x}+\frac{x}{2}-\frac{2x}{3}$\end{enumerate}
\label{m39392*uid56}\item Show that ${\left(2x-1\right)}^{2}-{\left(x-3\right)}^{2}$ can be simplified to $\left(x+2\right)\left(3x-4\right)$
\newline
\newline
\label{m39392*uid57}\item What must be added to ${x}^{2}-x+4$\hspace{1ex}to make it equal to ${\left(x+2\right)}^{2}$
\newline
\newline
\end{enumerate}
  \label{m39392**end}
  \label{d4e6ddcad4e2d9e383c4732da6858c66**end}
\practiceinfo
\par 
 \par \begin{tabular}[h]{cccccc}
 (1.) 00j9&  (2.) 00ja&  (3.) 00jb&  (4.) 00jc&  (5.) 00jd&  (6.) 00je& \end{tabular}
